% concatenated from Dunford-Pettis_Multilinear_Operators_and_their_variations__A_revisit_to_the_classic_concepts_of_Operator_Ideals_.tex
\documentclass[12 pt]{article}%
\usepackage{amsmath, amsfonts, amsthm, color,latexsym}
\usepackage{amsmath}
\usepackage{amsfonts}
\usepackage{amssymb}
\usepackage{color, soul}
%\usepackage[brazil]{babel}              % Determina a língua na qual o documento virá escrito.
%\usepackage[utf8]{inputenc}
%\usepackage{graphicx}
\usepackage{enumerate}
\setcounter{MaxMatrixCols}{30}
%TCIDATA{OutputFilter=latex2.dll}
%TCIDATA{Version=5.50.0.2890}
%TCIDATA{LastRevised=Sunday, March 28, 2010 09:24:22}
%TCIDATA{<META NAME="GraphicsSave" CONTENT="32">}
%TCIDATA{<META NAME="SaveForMode" CONTENT="1">}
%TCIDATA{BibliographyScheme=Manual}
%BeginMSIPreambleData
\providecommand{\U}[1]{\protect\rule{.1in}{.1in}}
%EndMSIPreambleData
\newtheorem{theorem}{Theorem}[section]
\newtheorem{proposition}[theorem]{Proposition}
\newtheorem{corollary}[theorem]{Corollary}
\newtheorem{example}[theorem]{Example}
\newtheorem{examples}[theorem]{Examples}
\newtheorem{remark}[theorem]{Remark}
\newtheorem{remarks}[theorem]{Remarks}
\newtheorem{lemma}[theorem]{Lemma}
\newtheorem{final remark}[theorem]{Final Remark}
\newtheorem{definition}[theorem]{Definition}
\textwidth=16.1cm
\textheight=23cm
\hoffset=-15mm
\voffset=-20mm
\def\blacksquare{\hbox{\vrule width 4pt height 4pt depth 0pt}}
\def\cqd{\ \ \ \hbox{}\nolinebreak\hfill $\blacksquare \  \  \
	\  $ \par{}\medskip}

\newcommand{\corr}[2]{\textcolor{blue}{\st{#1} \textcolor{red} {#2}}}
\allowdisplaybreaks[4]

\begin{document}
	
	%\title{\sc Weak and Weak* Dunford Pettis Multilinear Operators}
	\title{\sc Dunford-Pettis Multilinear Operators and their variations: A revisit to the classic concepts of Operator Ideals.}
	\author{Joilson Ribeiro,%\thanks{joilsonor@ufba.br}
		~ Fabr\'icio Santos \thanks{
			\newline\indent2010 Mathematics Subject
			Classification: 46B45, 46G25, 47H60, 47L22.\newline\indent Key words: Banach sequence spaces, Multipolynomias, Absolutely summing multipolynomials.}}
	\date{}
	\maketitle
	
	
	\begin{abstract}
		In this paper, we address broader concepts of Dunford-Pettis operators, presenting new classes and results that relate this class to others already well studied in the literature, as well as an approach outside the origin. We also investigate inclusion results and conditions for the coincidence of this class with previously studied classes and with new classes that arise throughout this work.
	\end{abstract}
	
	\section{Introduction and background}
	
	
	The class of Dunford–Pettis linear operators has played a crucial role in the theory of Banach spaces in general, as highlighted in \cite{DU77}. Several researchers have investigated this topic, mainly seeking correlations among the classes of compact linear operators, Dunford–Pettis operators, weakly Dunford–Pettis operators, and weakly* Dunford–Pettis operators, as can be seen in \cite{AB82, BHJ22, BKAB24}.
	
	A first multilinear approach to Dunford-Pettis operators between Banach spaces was proposed by Ribeiro, Santos, and Torres in \cite{RST22}, where it was shown that this class is an ideal of multilinear mappings but not a hyper-ideal, a concept introduced by Botelho and Torres in \cite{BT15}. In that work, Ribeiro, Santos, and Torres focused on verifying that this class does not satisfy the coherence concept introduced by Pellegrino and Ribeiro in \cite{PR14}. However, they did not explore other generalizations to the multilinear setting to study inclusion and coincidence results. It is worth noting that the case we will call Dunford-Pettis multilinear operators at every point already appears in the literature and can be found in \cite{AD02, BR98}.
	 
	In this paper, in addition to introducing new classes that generalize the class of Dunford-Pettis linear operators, we also revisit concepts such as multi-ideals, hyper-ideals, and coherence and compatibility to formalize these new structures into properties considered fundamental for their development, such as closure under composition with linear operators and the product of linear functionals and multilinear forms. We also present coincidence conditions for these classes and seek characterizations that make their use smoother.
	

	Let $E$ and $F$ be Banach spaces. An operator $T\colon E \rightarrow F$ is called a Dunford–Pettis operator, and is denoted by $\mathcal{DP}$, if $T$ carries weakly convergent sequences to norm convergent sequences. In other words, if $x_j \to 0$ as $j \to \infty$ in $\sigma(E,E')$, then $\|T(x_j)\|_F \to 0$ as $j \to \infty$.
	
	Alternatively, a (bounded linear) operator $T\colon E \rightarrow F$ is a Dunford–Pettis operator if and only if $T$ carries relatively weakly compact sets to norm totally bounded sets. A Banach space $E$ is said to have the Schur property if $Id_E\colon E \rightarrow E$, given by $Id_E(x)=x$, is a Dunford–Pettis operator. The Banach space $\ell_1$ has the Schur property (\cite[Theorem 2.5.18]{MN91}).
	
	We can also define the class of weakly Dunford–Pettis linear operators as follows: Let $E$ and $F$ be Banach spaces. An operator $T\colon E \rightarrow F$ is said to be weakly Dunford–Pettis if whenever $x_j \to 0$ in $\sigma(E,E')$ and $f_j \to 0$ in $F'$, then $f_j(T(x_j)) \to 0$. We denote this class by $\mathcal{WDP}$ (\cite{AB85}).
	
	From now on, the letters $E,E_{1},\dots,E_{m},F,G,H$ will denote Banach spaces over the same scalar field $\mathbb{K}=\mathbb{R}$ or $\mathbb{C}$, and $E'$ stands for the topological dual of $E$. The symbol $B_E$ denotes the closed unit ball of $E$. We use $\mathrm{BAN}$ to denote the class of all Banach spaces over $\mathbb{K}$.
	
	By $\mathcal{L}(E_1, \ldots, E_n;F)$ we denote the Banach space of continuous $n$-linear operators from $E_1 \times \cdots \times E_n$ to $F$, endowed with the usual uniform norm $\|\cdot\|$. In the linear case we write $\mathcal{L}(E;F)$. If $E_1 = \cdots = E_n = E$, we write $\mathcal{L}(^nE;F)$. If $F = \mathbb{K}$, we write $\mathcal{L}(E_1, \ldots, E_n)$ and $\mathcal{L}(^nE)$.
	
	A mapping $P \colon E \rightarrow F$ is said to be an {\it $m$-homogeneous polynomial} if there exists $A \in \mathcal{L}(^mE; F)$ such that $P(x) = A(x)^m$ for every $x \in E$, where $A(x)^m := A\left(x,\overset{m}{\dots},x\right)$. In this case we write $P = \widehat{A}$. By $\mathcal{P}(^mE;F)$ we denote the Banach space of all continuous $m$-homogeneous polynomials from $E$ to $F$, endowed with the usual supremum norm $\|\cdot\|$. By $\check{P}$ we denote the unique symmetric continuous $m$-linear operator associated with $P$.
	
	We write $x_j \overset{w}{\to} x$ in $E$ and $f_j \overset{w^*}{\to} f$ in $E'$ to mean that $x_j \to x$ in $\sigma(E,E')$ and $f_j \to f$ in $\sigma(E', E)$, as $j \to \infty$.
	
%	\begin{definition}
%		Let $m \in \mathbb{N}$. A Banach ideal of multilinear applications is a pair $\left(\mathcal{M}_m, \|\cdot \|_{\mathcal{M}_m} \right)$ where $\mathcal{M}_m$ is a  subclass of the class of all multilinear operators between Banach spaces and
%		\begin{equation*}
%			\|\cdot \|_{\mathcal{M}_m}\colon \mathcal{M}_m \longrightarrow \mathbb{R}
%		\end{equation*}
%		is a function such that, for all Banach spaces $E_1,\dots, E_m, F$, the component
%		
%		\begin{equation*}
%			\mathcal{M}_m(E_1,\dots, E_m; F) := \mathcal{L}(E_1,\dots, E_m; F) \cap \mathcal{M}_m
%		\end{equation*}
%		is a subspace of $\mathcal{L}(E_1,\dots, E_m; F)$ on which $\|\cdot \|_{\mathcal{M}_m}$ is a complete norm and
%		
%		\begin{enumerate}
%			\item The space of the multilinear operators of finite type is contained in $\mathcal{M}_m(E_1,\dots, E_m; F)$
%			
%			\item The application $I_m : \mathbb{K}^m \rightarrow \mathbb{K}$ given by $I_m(\lambda_1,\dots, \lambda_m) = \lambda_1 \cdots \lambda_m$ belongs  to $\mathcal{M}_m(\mathbb{K}^m; \mathbb{K})$ and
%			
%			\begin{equation*}
%				\|I_m\|_{\mathcal{M}_m} = 1.
%			\end{equation*}
%			
%			\item $($Multi-ideal property$)$ If $T \in \mathcal{M}_m(E_1,\dots, E_m; F), u_i \in \mathcal{L}(G_i; E_i), i=1,\dots, m$ and $t \in \mathcal{L}(F; H)$ then  $t\circ T \circ (u_1,\dots, u_m) \in \mathcal{M}_m(G_1,\dots, G_m; H)$ and
%			
%			\begin{equation*}
%				\|t\circ T \circ (u_1,\dots, u_m) \|_{\mathcal{M}_m} \le \|t\| \left\|T \right\|_{\mathcal{M}_m} \|u_1\|\cdots \|u_m\|.
%			\end{equation*}
%		\end{enumerate}
%	\end{definition}
%
%Analogously, we can define the homogeneous polynomial ideal $\mathcal{Q}$.  An ideal of homogeneous polynomials $\mathcal{Q}$ is a subclass of the class $\mathcal{P}={\textstyle\bigcup\limits_{n=1}^{\infty}} \mathcal{P}_{n}$ of all continuous homogeneous polynomials between Banach spaces, such that, for a positive integer $n$, Banach spaces $E$ and $F$, the components
%\[
%\mathcal{Q}_{n}(E; F):=\mathcal{P}_{n}(^nE; F)\cap\mathcal{P}%
%\]
%satisfy:
%
%(Ma) $\mathcal{Q}_{n}(^nE; F)$ is a linear subspace of $\mathcal{P}_{n}(^nE; F)$, which contains the $n$-homogeneous finite type polynomials, where an $n$-homogeneous polynomials is said finite type when
%\begin{equation*}
%	P(x) = \sum_{i=1}^k[\varphi_i(x)]^nb_i
%\end{equation*}
%with $k \in \mathbb{N}$ and $\varphi_i \in E'$, $b_i \in F$, $i = 1,\dots, k$.
%
%(Mb) If $P\in\mathcal{Q}_{n}(^nE;F)$, $u \in\mathcal{L}(G;E)$ and $v\in\mathcal{L}(F;H)$, then
%\[
%v\circ P\circ u\in\mathcal{Q}_{n}(^nG; H).
%\]
%
%Moreover, $\mathcal{Q}$ is a (quasi-) normed multi-ideal if there is a function $\Vert\cdot\Vert_{\mathcal{Q}}\colon\mathcal{Q}\longrightarrow\lbrack0,\infty)$ satisfying
%
%(M1) $\Vert\cdot\Vert_{\mathcal{Q}}$ restricted to $\mathcal{Q}_{n}(^nE; F)$ is a (quasi-) norm, for all Banach spaces $E$ and $F.$
%
%(M2) $\Vert P_{n}\colon\mathbb{K}^{n}\longrightarrow\mathbb{K}:P_{n}(\lambda_{1},\ldots,\lambda_{n})=\lambda^n\Vert_{\mathcal{Q}}=1$ for all $n$,
%
%(M3) If $P\in\mathcal{Q}_{n}(^nE; F)$, $u \in\mathcal{L}(G; E)$  and $v\in\mathcal{L}(F;H)$, then
%\[
%\Vert v\circ P\circ u \Vert_{\mathcal{Q}}\leq\Vert
%v\Vert\Vert P\Vert_{\mathcal{Q}}\Vert u\Vert^n.
%\]
%When all of the components $\mathcal{Q}_{n}(^nE; F)$ are complete under this (quasi-) norm, $\mathcal{Q}$ is called the (quasi-) Banach homogeneous polynomials ideal. For a fixed homogeneous polynomials ideal $\mathcal{Q}$ and a positive integer $n$, the class
%\[
%\mathcal{Q}_{n}:=\cup_{E,F}\mathcal{Q}_{n}\left(^nE; F\right)
%\]
%is called  of $n$-homogeneous polynomials ideal. For more details, see \cite{FG03}.


\section{Dunford-Pettis Multilinear Operators and Homogeneous Polynomials Everywhere}\label{DPMOHPEv}

In this section we present the concept of Dunford–Pettis multilinear operators and Dunford–Pettis homogeneous polynomials at every point. Recall that the concept of Dunford–Pettis multilinear operators and homogeneous polynomials (at the origin) was introduced by Ribeiro, Santos, and Torres in \cite{RST22}. In that paper it was shown that the coherence condition (CH1) fails, and that this class does not generate a hyper-ideal of multilinear mappings or homogeneous polynomials, concepts introduced in \cite{BT15} and \cite{BT18}. We begin by defining these notions.

\begin{definition}\label{DPML}
	Let $m \in \mathbb{N}$, let $E_1,\dots, E_m$ and $F$ be Banach spaces, and let $T \in \mathcal{L}(E_1,\dots, E_m; F)$. We say that $T$ is a Dunford–Pettis $m$-linear operator at $(a_1,\dots, a_m) \in E_1 \times \cdots \times E_m$ if, for each $i=1,\dots,m$, every sequence $(x_j^{(i)})_{j=1}^\infty \subset E_i$ satisfying
	\[
	x_j^{(i)} \overset{w}{\longrightarrow} a_i
	\]
	also satisfies
	\[
	T(x_j^{(1)},\dots,x_j^{(m)}) \longrightarrow T(a_1,\dots,a_m)
	\]
	in norm as $j \to \infty$.
\end{definition}

If $T$ is a Dunford–Pettis multilinear operator at every point of $E_1 \times \cdots \times E_m$, we say that $T$ is Dunford–Pettis at every point. We denote by $\mathcal{L_{DP}^{\mathrm{ev}}}$ the class of all such operators. Analogously, we have the following definition.

\begin{definition}\label{DPPH}
	Let $m \in \mathbb{N}$, let $E$ and $F$ be Banach spaces, and let $P \in \mathcal{P}(^mE; F)$. We say that $P$ is a Dunford–Pettis $m$-homogeneous polynomial at $a \in E$ if, for every sequence $(x_j)_{j=1}^{\infty} \subset E$ with
	\begin{equation*}
		x_j \overset{w}{\longrightarrow} a,
	\end{equation*}
	we have
	\begin{equation*}
		P(x_j) \longrightarrow P(a)
	\end{equation*}
	in norm as $j \to \infty$.
\end{definition}

If $P$ is a Dunford–Pettis homogeneous polynomial at every point of $E$, we say that $P$ is Dunford–Pettis at every point. We denote by $\mathcal{P_{DP}^{\mathrm{ev}}}$ the class of all such polynomials.

\begin{proposition}
	Let $E_1,\dots, E_m$ and $F$ be Banach spaces and $T \in \mathcal{L}(E_1,\dots, E_m; F)$. Then, $T \in \mathcal{L_{DP}^{\text{ev}}}(E_1,\dots, E_m; F)$ if, and only if, for every $(a_1,\dots, a_m) \in E_1\times\cdots\times E_m$ and every weakly null sequences $(x_j^{(i)})_{j=1}^{\infty}$ in $E_i$, $i=1,\dots, m$, we have
	\begin{equation*}
		\|T(x_j^{(1)} + a_1,\dots, x_j^{(m)} + a_m) - T(a_1,\dots, a_m)\|\longrightarrow 0.
	\end{equation*}
\end{proposition}

\begin{proof}
	Suppose first that $T \in \mathcal{L_{DP}^{\text{ev}}}(E_1,\dots, E_m; F)$. 
	
	Let $(a_1,\dots,a_m)\in E_1\times\cdots\times E_m$ and let $(x_j^{(i)})$ be weakly null sequences in $E_i$. Then
	\[
	x_j^{(i)} + a_i \overset{w}{\longrightarrow} a_i,
	\]
	for all $i=1,\dots,m$. Since $T$ is Dunford-Pettis at every point, it follows that
	\[
	\|T(x_j^{(1)} + a_1,\dots, x_j^{(m)} + a_m) - T(a_1,\dots, a_m)\| \to 0.
	\]
	
	Conversely, suppose the condition holds. Let $(x_j^{(i)}) \subset E_i$ be such that $x_j^{(i)} \overset{w}{\longrightarrow} a_i$. Define
	\[
	y_j^{(i)} = x_j^{(i)} - a_i,
	\]
	so that $y_j^{(i)} \overset{w}{\longrightarrow} 0$. Then
	\begin{align*}
		&\|T(x_j^{(1)},\dots, x_j^{(m)}) - T(a_1,\dots, a_m)\| \\
		&= \|T(y_j^{(1)} + a_1,\dots, y_j^{(m)} + a_m) - T(a_1,\dots, a_m)\|.
	\end{align*}
	By hypothesis, the right-hand side converges to $0$, and hence
	\[
	T(x_j^{(1)},\dots, x_j^{(m)}) \to T(a_1,\dots, a_m).
	\]
	Therefore, $T \in \mathcal{L_{DP}^{\text{ev}}}(E_1,\dots, E_m; F)$.
\end{proof}

\begin{remark}\label{DPevcontDP}
	\begin{itemize}
		\item[(1)] The concepts in Definitions \ref{DPML} and \ref{DPPH} are not new (\cite{BT15, BT18}). They coincide with the classes of weakly sequentially continuous multilinear operators and weakly sequentially continuous homogeneous polynomials, respectively. We adopt this terminology for alignment with the class of Dunford–Pettis multilinear mappings.
		
		\item[(2)] It is immediate that $\mathcal{L_{DP}^{\mathrm{ev}}} \subset \mathcal{L_{DP}}$ and $\mathcal{P_{DP}^{\mathrm{ev}}} \subset \mathcal{P_{DP}}$.
	\end{itemize}
\end{remark}

%\begin{remark}\label{DPevcontDP}
%	É imediato perceber que $\mathcal{L_{DP}^{\text{ev}}} \subset \mathcal{L_{DP}}$ e também que $\mathcal{P_{DP}^{\text{ev}}} \subset \mathcal{P_{DP}}$.
%\end{remark}

A natural question at this point is whether $\mathcal{L_{DP}^{\mathrm{ev}}}$ is a Banach multi-ideal. To address this question, we first show that $\mathcal{L_{DP}^{\mathrm{ev}}}(E_1,\dots,E_m;F)$ is a Banach space for arbitrary Banach spaces $E_1,\dots,E_m$ and $F$.

\begin{proposition}
	Let $m \in \mathbb{N}$, let $E_1,\dots,E_m$ and $F$ be Banach spaces. Then $\mathcal{L_{DP}^{\mathrm{ev}}}(E_1,\dots,E_m;F)$ is a closed subset of $\mathcal{L_{DP}}(E_1,\dots,E_m;F)$.
\end{proposition}
\begin{proof}
	Let $(T_j)_{j=1}^{\infty} \subset \mathcal{L_{DP}^{\mathrm{ev}}}(E_1,\dots,E_m;F)$ with $T_j \to T$ in norm. Let $(x_j^{(i)})_{j=1}^{\infty} \subset E_i$ be such that $x_j^{(i)} \overset{w}{\to} x_i$, $i=1,\dots,m$. Then $(x_j^{(i)})$ are bounded, so
	\[
	\|(x_j^{(1)},\dots,x_j^{(m)})\| := \|x_j^{(1)}\|\cdots\|x_j^{(m)}\| \le C
	\]
	for some $C>0$ and all $j$.
	
	Given $\varepsilon>0$, choose $k$ such that $\|T-T_k\|<\varepsilon$. Since $T_k \in \mathcal{L_{DP}^{\mathrm{ev}}}$,
	\[
	T_k(x_j^{(1)},\dots,x_j^{(m)}) \to T_k(x_1,\dots,x_m).
	\]
	Hence,
	\begin{align*}
		&\|T(x_j^{(1)},\dots,x_j^{(m)}) - T(x_1,\dots,x_m)\| \\
		&\le C\|T-T_k\| + \|T_k(x_j^{(1)},\dots,x_j^{(m)}) - T_k(x_1,\dots,x_m)\| + \|T-T_k\|\|(x_1,\dots,x_m)\|,
	\end{align*}
	which tends to $0$ as $j \to \infty$. Therefore,
	\[
	T(x_j^{(1)},\dots,x_j^{(m)}) \to T(x_1,\dots,x_m),
	\]
	and $T \in \mathcal{L_{DP}^{\mathrm{ev}}}(E_1,\dots,E_m;F)$.
\end{proof}

%\begin{proof}
%	Let $(T_k)_{k=1}^{\infty} \subset \mathcal{L_{DP}^{\mathrm{ev}}}(E_1,\dots,E_m;F)$ such that $T_k \to T$ in norm. Let $(x_j^{(i)})_{j=1}^{\infty} \subset E_i$ with $x_j^{(i)} \overset{w}{\to} x_i$, $i=1,\dots,m$. Then each $(x_j^{(i)})$ is bounded, so there exist $K_i>0$ such that
%	\[
%	\|x_j^{(i)}\| \le K_i \quad \text{for all } j \text{ and } i=1,\dots,m.
%	\]
%	Thus, defining
%	\[
%	\|(x_j^{(1)},\dots,x_j^{(m)})\| := \|x_j^{(1)}\|\cdots\|x_j^{(m)}\|,
%	\]
%	we have
%	\[
%	\|(x_j^{(1)},\dots,x_j^{(m)})\| \le C := K_1\cdots K_m \quad \text{for all } j.
%	\]
%	
%	Let $\varepsilon > 0$. Choose $k_0$ such that
%	\[
%	\|T - T_{k_0}\| < \frac{\varepsilon}{2(C + \|(x_1,\dots,x_m)\|)}.
%	\]
%	Since $T_{k_0} \in \mathcal{L_{DP}^{\mathrm{ev}}}$, there exists $j_0$ such that for all $j \ge j_0$,
%	\[
%	\|T_{k_0}(x_j^{(1)},\dots,x_j^{(m)}) - T_{k_0}(x_1,\dots,x_m)\| < \frac{\varepsilon}{2}.
%	\]
%	
%	Then, for all $j \ge j_0$,
%	\begin{align*}
%		&\|T(x_j^{(1)},\dots,x_j^{(m)}) - T(x_1,\dots,x_m)\| \\
%		&\le \|T - T_{k_0}\|\,\|(x_j^{(1)},\dots,x_j^{(m)})\| \\
%		&\quad + \|T_{k_0}(x_j^{(1)},\dots,x_j^{(m)}) - T_{k_0}(x_1,\dots,x_m)\| \\
%		&\quad + \|T - T_{k_0}\|\,\|(x_1,\dots,x_m)\| \\
%		&< \frac{\varepsilon}{2} + \frac{\varepsilon}{2} = \varepsilon.
%	\end{align*}
%	
%	Therefore,
%	\[
%	T(x_j^{(1)},\dots,x_j^{(m)}) \to T(x_1,\dots,x_m),
%	\]
%	and hence $T \in \mathcal{L_{DP}^{\mathrm{ev}}}(E_1,\dots,E_m;F)$.
%\end{proof}





%\begin{proposition}
%	Let $m \in \mathbb{N}$, let $E_1,\dots,E_m$ and $F$ be Banach spaces. Then $\mathcal{L_{DP}^{\mathrm{ev}}}(E_1,\dots,E_m;F)$ is a closed subset of $\mathcal{L_{DP}}(E_1,\dots,E_m;F)$.
%\end{proposition}
%
%\begin{proof}
%	Let $(T_j)_{j=1}^{\infty}$ be a sequence in $\mathcal{L_{DP}^{\mathrm{ev}}}(E_1,\dots,E_m;F)$ converging to $T \in \mathcal{L}(E_1,\dots,E_m;F)$. Let $(x_n^{(i)})_{n=1}^{\infty} \subset E_i$ be sequences such that $x_n^{(i)} \overset{w}{\to} x_i$ in $E_i$, for each $i=1,\dots,m$. As
%	\begin{align*}
%		&\|T(x_n^{(1)},\dots, x_n^{(m)}) - T(x_1,\dots, x_m)\|\le \\ 
%		%		&\le \|T(x_n^{(1)},\dots, x_n^{(m)}) - T_j(x_n^{(1)},\dots, x_n^{(m)})\| +\\ 
%		%		&+ \|T_j(x_n^{(1)},\dots, x_n^{(m)}) - T_j(x_{1},\dots, x_{m})\| + \|T_j(x_{1},\dots, x_{m}) - T(x_{1},\dots, x_{m})\|\\
%		&\|T_j - T\| \|(x_n^{(1)},\dots, x_n^{(m)})\| + \|T_j(x_n^{(1)},\dots, x_n^{(m)}) - T_j(x_{1},\dots, x_{m})\| + \|T_j - T\|\|(x_{1},\dots, x_{m})\|
%	\end{align*}
%	and $T_j \in \mathcal{L_{DP}^{\text{ev}}}(E_1,\dots, E_m; F)$ for all $j \in \mathbb{N}$, then
%	\begin{equation*}
%		T(x_n^{(1)},\dots, x_n^{(m)}) \overset{\| \cdot \|}{\longrightarrow} T(x_1,\dots, x_m)
%	\end{equation*}
%	as $n \rightarrow \infty$. So, $T \in \mathcal{L_{DP}^{\text{ev}}}(E_1,\dots, E_m; F)$.
%\end{proof}

As an immediate consequence of this proposition, and since $\mathcal{L_{DP}}(E_1,\dots,E_m;F)$ is a Banach space, we obtain the following result.

\begin{corollary}
	Let $m \in \mathbb{N}$, let $E_1,\dots,E_m$ and $F$ be Banach spaces. Then $\mathcal{L_{DP}^{\mathrm{ev}}}(E_1,\dots,E_m;F)$ is a Banach space.
\end{corollary}

It is straightforward to verify that the pairs $(\mathcal{L_{DP}^{\mathrm{ev}}}, \|\cdot\|)$ and $(\mathcal{P_{DP}^{\mathrm{ev}}}, \|\cdot\|)$ form a Banach multi-ideal and a Banach ideal of homogeneous polynomials, respectively (see \cite{FG03, P83} for details). These structures generalize the Banach ideal of linear operators $(\mathcal{L_{DP}}, \|\cdot\|)$. On the other hand, $(\mathcal{L_{DP}^{\mathrm{ev}}}, \|\cdot\|)$ is not a hyper-ideal of multilinear mappings, as shown in Examples 4.4 and 3.11 of \cite{BT15} and \cite{BT18}, respectively.

We next show that the sequence
\begin{equation*}
	\left(\mathcal{L_{DP}^{\mathrm{ev,n}}}, \mathcal{P_{DP}^{\mathrm{ev,n}}} \right)_{n=1}^{\infty}
\end{equation*}
is coherent and compatible with $(\mathcal{L_{DP}}, \|\cdot\|)$, according to \cite{PR14}. However, it is not strongly coherent (see \cite{RST22}). We begin by establishing the following result.

Before proving the next result, we recall a standard notion in multilinear theory. Let $S_m$ denote the group of permutations of $\{1,\dots,m\}$. Given $T \in \mathcal{L}(^mE;F)$ and $\sigma \in S_m$, define $T_{\sigma}$ and $T_S$ in $\mathcal{L}(^mE;F)$ by
\begin{align*}
	T_{\sigma}(x_1,\dots,x_m) &= T(x_{\sigma(1)},\dots,x_{\sigma(m)}),\\
	T_S(x_1,\dots,x_m) &= \frac{1}{m!}\sum_{\sigma \in S_m} T_{\sigma}(x_1,\dots,x_m).
\end{align*}

A subclass $\mathcal{G}$ of the class of continuous multilinear operators between Banach spaces is said to be symmetric if $T_S \in \mathcal{G}(^mE;F)$ whenever $T \in \mathcal{G}(^mE;F)$. If $\mathcal{G}$ is endowed with a function $\|\cdot\|_{\mathcal{G}}\colon \mathcal{G} \to [0,\infty)$, then $\mathcal{G}$ is said to be strongly symmetric if $T_{\sigma} \in \mathcal{G}$ and $\|T_{\sigma}\|_{\mathcal{G}} = \|T\|_{\mathcal{G}}$ for all $T \in \mathcal{G}(^mE;F)$ and $\sigma \in S_m$ (see \cite{FG03}).

\begin{proposition}\label{DPS}
	The subclass $\mathcal{L_{DP}^{\mathrm{ev}}}$ of the class of continuous multilinear operators between Banach spaces is symmetric.
\end{proposition}
 
\begin{proof}
	Let $E_1,\dots,E_m$ and $F$ be Banach spaces, let $\sigma \in S_m$, and let $T \in \mathcal{L_{DP}^{\mathrm{ev}}}(E_1,\dots,E_m;F)$. Let $(x_j^{(i)})_{j=1}^{\infty} \subset E_i$ be sequences such that
	\begin{equation*}
		x_j^{(i)} \overset{w}{\longrightarrow} x_i \in E_i, \quad i=1,\dots,m.
	\end{equation*}
	Then
	\begin{equation*}
		T(x_j^{(1)},\dots,x_j^{(m)}) \longrightarrow T(x_1,\dots,x_m)
	\end{equation*}
	in norm as $j \to \infty$.
	
	Since $x_j^{(\sigma(i))} \overset{w}{\to} x_{\sigma(i)}$ for $i=1,\dots,m$, it follows that
	\begin{equation*}
		\|T_{\sigma}(x_j^{(1)},\dots,x_j^{(m)}) - T_{\sigma}(x_1,\dots,x_m)\|
		= \|T(x_j^{(\sigma(1))},\dots,x_j^{(\sigma(m))}) - T(x_{\sigma(1)},\dots,x_{\sigma(m)})\| \to 0
	\end{equation*}
	as $j \to \infty$. Hence,
	\begin{equation*}
		\|T_S(x_j^{(1)},\dots,x_j^{(m)}) - T_S(x_1,\dots,x_m)\| \to 0
	\end{equation*}
	as $j \to \infty$. Therefore, $T_S \in \mathcal{L_{DP}^{\mathrm{ev}}}(E_1,\dots,E_m;F)$.
\end{proof}

\begin{theorem}\label{DPcoercomp}
 The sequence $\left(\mathcal{L_{DP}^{\mathrm{ev,n}}}, \mathcal{P_{DP}^{\mathrm{ev,n}}} \right)_{n=1}^{\infty}$ is coherent and compatible with $\mathcal{DP}$.
\end{theorem}

\begin{proof}
		
Let $m \in \mathbb{N}$, and let $E$ and $F$ be Banach spaces. We begin by showing that, for any $k \in \{1,\dots,m\}$, $P \in \mathcal{P_{DP}^{\mathrm{ev}}}(^kE;F)$ if and only if $\check{P} \in \mathcal{L_{DP}^{\mathrm{ev}}}(^kE;F)$ (see Properties (CH5) and (CP5) in \cite[Definitions 2.1 and 2.2]{PR14}).

Suppose that $\check{P} \in \mathcal{L_{DP}^{\mathrm{ev}}}(^kE;F)$. Let $(x_j)_{j=1}^{\infty} \subset E$ be such that $x_j \overset{w}{\to} x \in E$. Then
\begin{equation*}
	P(x_j) = \check{P}(x_j,\dots,x_j) \longrightarrow \check{P}(x,\dots,x) = P(x)
\end{equation*}
in norm, and hence $P \in \mathcal{P_{DP}^{\mathrm{ev}}}(^kE;F)$.

Conversely, suppose that $P \in \mathcal{P_{DP}^{\mathrm{ev}}}(^kE;F)$. Let $(x_j^{(i)})_{j=1}^{\infty} \subset E$ be such that $x_j^{(i)} \overset{w}{\to} x_i \in E$, for all $i=1,\dots,k$. Since $\epsilon_i = \pm 1$, we have
\begin{equation*}
	\epsilon_1 x_j^{(1)} + \cdots + \epsilon_k x_j^{(k)} \overset{w}{\longrightarrow} \epsilon_1 x_1 + \cdots + \epsilon_k x_k.
\end{equation*}
Hence,
\begin{align*}
	&\left\|\check{P}(x_j^{(1)},\dots,x_j^{(k)}) - \check{P}(x_1,\dots,x_k)\right\| \\
	&= \frac{1}{2^k k!} \left\| P\!\left(\sum_{i=1}^k \epsilon_i x_j^{(i)}\right) - P\!\left(\sum_{i=1}^k \epsilon_i x_i\right) \right\| \longrightarrow 0
\end{align*}
as $j \to \infty$. Thus, $\check{P} \in \mathcal{L_{DP}^{\mathrm{ev}}}(^kE;F)$.

Next, we show that for any $T \in \mathcal{L_{DP}^{\mathrm{ev}}}(E_1,\dots,E_m;F)$ and $a_i \in E_i$, we have $T_{a_i} \in \mathcal{L_{DP}^{\mathrm{ev}}}(E_1,\dots,E_{i-1},E_{i+1},\dots,E_m;F)$ and
\begin{equation*}
	\|T_{a_i}\| \le \|T\| \|a_i\|,
\end{equation*}
for all $i=1,\dots,m$ (Property (CH1) in \cite[Definition 2.1]{PR14}). Without loss of generality, consider $i=1$. Let $(x_j^{(i)})_{j=1}^{\infty} \subset E_i$ be such that $x_j^{(i)} \overset{w}{\to} x_i$ for $i=2,\dots,m$. Then
\begin{equation*}
	\|T_{a_1}(x_j^{(2)},\dots,x_j^{(m)}) - T_{a_1}(x_2,\dots,x_m)\|
	= \|T(a_1,x_j^{(2)},\dots,x_j^{(m)}) - T(a_1,x_2,\dots,x_m)\| \to 0,
\end{equation*}
since the constant sequence $(x_j^{(1)})_{j=1}^{\infty}=a_1$ converges weakly to $a_1$. The norm inequality is immediate.

Finally, let $T \in \mathcal{L_{DP}^{\mathrm{ev}}}(E_1,\dots,E_m;F)$ and $\gamma \in E_{m+1}'$. We show that $\gamma T \in \mathcal{L_{DP}^{\mathrm{ev}}}(E_1,\dots,E_{m+1};F)$. Let $(x_j^{(i)})_{j=1}^{\infty} \subset E_i$ be such that $x_j^{(i)} \overset{w}{\to} x_i \in E_i$ for all $i=1,\dots,m+1$. Then
\begin{align*}
	&\gamma T(x_j^{(1)},\dots,x_j^{(m+1)}) - \gamma T(x_1,\dots,x_{m+1}) \\
	&= \gamma(x_j^{(m+1)} - x_{m+1})\, T(x_j^{(1)},\dots,x_j^{(m)}) \\
	&\quad + \gamma(x_{m+1})\big(T(x_j^{(1)},\dots,x_j^{(m)}) - T(x_1,\dots,x_m)\big).
\end{align*}
Since $(T(x_j^{(1)},\dots,x_j^{(m)}))$ is bounded and $x_j^{(m+1)} \overset{w}{\to} x_{m+1}$, it follows that
\begin{align*}
	&\|\gamma T(x_j^{(1)},\dots,x_j^{(m+1)}) - \gamma T(x_1,\dots,x_{m+1})\| \\
	&\le |\gamma(x_j^{(m+1)} - x_{m+1})|\, \|T(x_j^{(1)},\dots,x_j^{(m)})\| \\
	&\quad + |\gamma(x_{m+1})|\, \|T(x_j^{(1)},\dots,x_j^{(m)}) - T(x_1,\dots,x_m)\| \to 0
\end{align*}
as $j \to \infty$. The norm estimate is immediate.

The conclusion follows from \cite[Remark 2.3(3) and Proposition 2.4]{RST22} together with Proposition \ref{DPS}.
	
\end{proof}

\begin{remark}
	We present an example showing that this class is not strongly coherent. Let $E_1,\dots,E_m$ and $F$ be Banach spaces, let $\varphi_i \in E_i'$, $i=1,\dots,m$, and let $y \in F$ be nonzero. Choose $(a_1,\dots,a_m) \in E_1 \times \cdots \times E_m$ such that $a_i \notin \ker \varphi_i$ for all $i$.
	
	Define $T \colon E_1 \times \cdots \times E_m \to F$ by
	\begin{equation*}
		T(x_1,\dots,x_m) = \varphi_1(x_1)\cdots \varphi_m(x_m)\, y,
	\end{equation*}
	and let $Q \colon \ell_2 \times \ell_2 \to \mathbb{K}$ be given by
	\begin{equation*}
		Q\big((x_j)_{j=1}^{\infty}, (y_j)_{j=1}^{\infty}\big) = \sum_{j=1}^{\infty} x_j y_j.
	\end{equation*}
	
	Let $(x_j^{(i)})_{j=1}^{\infty} \subset E_i$ be such that $x_j^{(i)} \overset{w}{\to} a_i$ for $i=1,\dots,m$. Since $e_j \overset{w}{\to} 0$ in $\ell_2$, we have
	\begin{align*}
		QT(x_j^{(1)},\dots,x_j^{(m)},e_j,e_j)
		&= Q(e_j,e_j)\, T(x_j^{(1)},\dots,x_j^{(m)}) \\
		&= \varphi_1(x_j^{(1)})\cdots \varphi_m(x_j^{(m)})\, y \\
		&\longrightarrow \varphi_1(a_1)\cdots \varphi_m(a_m)\, y \neq 0.
	\end{align*}
	Thus,
	\begin{equation*}
		QT(x_j^{(1)},\dots,x_j^{(m)},e_j,e_j) \nrightarrow QT(a_1,\dots,a_m,0,0),
	\end{equation*}
	and therefore $QT \notin \mathcal{L_{DP}^{\mathrm{ev}}}(E_1,\dots,E_m,\ell_2,\ell_2;F)$, even though $T \in \mathcal{L_{DP}^{\mathrm{ev}}}(E_1,\dots,E_m;F)$.
\end{remark}

%\color{red}


%\begin{proposition}
%	Let $E_1,\dots, E_m$ and $F$ be Banach spaces, $(a_1,\dots, a_m) \in E_1\times\cdots\times E_m$ and $T \in \mathcal{L}(E_1,\dots, E_m; F)$. Then, $T \in \mathcal{L_{DP}^{\text{ev}}}(E_1,\dots, E_m; F)$ if, and only if, for any weakly null sequences $(x_j^{(i)})_{j=1}^{\infty}$ in $E_i$, for all $i = 1,\dots, m$, we have
%	\begin{equation*}
%		\|T(x_j^{(1)} + a_1,\dots, x_j^{(m)} + a_m) - T(a_1,\dots, a_m)\|\longrightarrow 0
%	\end{equation*}
%	as $j \rightarrow \infty$.
%\end{proposition}
%
%\begin{proof}
%	Suppose first that $T \in \mathcal{L_{DP}^{\text{ev}}}(E_1,\dots, E_m; F)$. 
%	We treat the case $m=3$, the general case being analogous.
%	
%	Let $(x_j^{(i)})_{j=1}^{\infty}$ be weakly null sequences in $E_i$ and fix $a_i \in E_i$, $i=1,2,3$. By multilinearity,
%	\begin{align*}
%		&T(x_j^{(1)}+a_1, x_j^{(2)}+a_2, x_j^{(3)}+a_3) - T(a_1,a_2,a_3) \\
%		&= T(x_j^{(1)}, x_j^{(2)}, x_j^{(3)}) 
%		+ T_{a_3}(x_j^{(1)}, x_j^{(2)}) 
%		+ T_{a_2}(x_j^{(1)}, x_j^{(3)}) \\
%		&\quad + T_{a_1}(x_j^{(2)}, x_j^{(3)}) 
%		+ T_{a_2,a_3}(x_j^{(1)}) 
%		+ T_{a_1,a_3}(x_j^{(2)}) 
%		+ T_{a_1,a_2}(x_j^{(3)}).
%	\end{align*}
%	
%	Since $T$ is Dunford-Pettis at every point, each term converges to $0$, hence
%	\[
%	\|T(x_j^{(1)}+a_1, x_j^{(2)}+a_2, x_j^{(3)}+a_3) - T(a_1,a_2,a_3)\| \to 0.
%	\]
%	
%	Conversely, suppose the condition holds. Let $x_j^{(i)} \overset{w}{\to} a_i$ and define $y_j^{(i)} = x_j^{(i)} - a_i$. Then $y_j^{(i)} \overset{w}{\to} 0$ and
%	\[
%	\|T(x_j^{(1)},\dots,x_j^{(m)}) - T(a_1,\dots,a_m)\|
%	= \|T(y_j^{(1)}+a_1,\dots,y_j^{(m)}+a_m) - T(a_1,\dots,a_m)\| \to 0.
%	\]
%	
%	Therefore, $T \in \mathcal{L_{DP}^{\text{ev}}}(E_1,\dots,E_m;F)$.
%\end{proof}


%\begin{lemma}[Characterization via weakly null perturbations]\label{DPchar}
%	Let $E_1,\dots, E_m$ and $F$ be Banach spaces and $T \in \mathcal{L}(E_1,\dots, E_m; F)$. The following are equivalent:
%	\begin{itemize}
%		\item[(i)] $T \in \mathcal{L_{DP}^{\mathrm{ev}}}(E_1,\dots, E_m; F)$;
%		
%		\item[(ii)] For every $(a_1,\dots, a_m) \in E_1 \times \cdots \times E_m$ and every family of weakly null sequences $(x_j^{(i)})_{j=1}^\infty \subset E_i$, $i=1,\dots,m$, we have
%		\begin{equation*}
%			\|T(x_j^{(1)}+a_1,\dots,x_j^{(m)}+a_m)-T(a_1,\dots,a_m)\|\longrightarrow 0.
%		\end{equation*}
%	\end{itemize}
%\end{lemma}
%
%\begin{proof}
%	$(i)\Rightarrow(ii)$:
%	Fix $(a_1,\dots,a_m)$ and weakly null sequences $(x_j^{(i)})$. Then
%	\[
%	x_j^{(i)}+a_i \overset{w}{\longrightarrow} a_i,
%	\]
%	so, by the definition of $\mathcal{L_{DP}^{\mathrm{ev}}}$,
%	\[
%	\|T(x_j^{(1)}+a_1,\dots,x_j^{(m)}+a_m)-T(a_1,\dots,a_m)\|\to 0.
%	\]
%	
%	$(ii)\Rightarrow(i)$:
%	Let $x_j^{(i)} \overset{w}{\longrightarrow} a_i$ and define $y_j^{(i)}=x_j^{(i)}-a_i$. Then $y_j^{(i)} \overset{w}{\longrightarrow} 0$ and
%	\[
%	T(x_j^{(1)},\dots,x_j^{(m)}) - T(a_1,\dots,a_m)
%	= T(y_j^{(1)}+a_1,\dots,y_j^{(m)}+a_m)-T(a_1,\dots,a_m).
%	\]
%	By hypothesis, the right-hand side converges to $0$, proving that $T \in \mathcal{L_{DP}^{\mathrm{ev}}}$.
%\end{proof}

%\color{black}

We denote by $\mathcal{L_{DP}^{\mathrm{(CH1)}}}$ the class of all multilinear Dunford–Pettis operators between Banach spaces that satisfy condition (CH1) (see \cite[Definition 3.2]{PR14} and \cite[Definition 2.1]{RST22}).

\begin{theorem}\label{DPCH1}
	Let $E_1,\dots,E_m$ and $F$ be Banach spaces. Then
	\begin{equation*}
		\mathcal{L_{DP}^{\mathrm{(CH1)}}} = \mathcal{L_{DP}^{\mathrm{ev}}}.
	\end{equation*}
\end{theorem}

\begin{proof}
	Let $E_1,\dots,E_m$ and $F$ be Banach spaces. It follows from Remark \ref{DPevcontDP} and Theorem \ref{DPcoercomp} that
	\begin{equation*}
		\mathcal{L_{DP}^{\mathrm{ev}}} \subset \mathcal{L_{DP}^{\mathrm{(CH1)}}}.
	\end{equation*}
	
	We now prove the reverse inclusion. Let $T \in \mathcal{L_{DP}^{\mathrm{(CH1)}}}(E_1,\dots,E_m;F)$. We show that $T \in \mathcal{L_{DP}^{\mathrm{ev}}}(E_1,\dots,E_m;F)$. Let $(x_j^{(i)})_{j=1}^{\infty} \subset E_i$ be such that
	\begin{equation*}
		x_j^{(i)} \overset{w}{\to} a_i \in E_i,
	\end{equation*}
	for all $i=1,\dots,m$. Then $(x_j^{(i)} - a_i) \overset{w}{\to} 0$ in $E_i$ for all $i=1,\dots,m$.
	
	For simplicity, we consider the case $m=3$; the general case follows analogously. We have
	\begin{align*}
		&T(x_j^{(1)},x_j^{(2)},x_j^{(3)}) - T(a_1,a_2,a_3) \\
		&= T(x_j^{(1)}-a_1, x_j^{(2)}-a_2, x_j^{(3)}-a_3) \\
		&\quad + T_{a_2,a_3}(x_j^{(1)}-a_1) + T_{a_1,a_3}(x_j^{(2)}-a_2) + T_{a_1,a_2}(x_j^{(3)}-a_3) \\
		&\quad + T_{a_3}(x_j^{(1)}-a_1, x_j^{(2)}-a_2) + T_{a_2}(x_j^{(1)}-a_1, x_j^{(3)}-a_3) + T_{a_1}(x_j^{(2)}-a_2, x_j^{(3)}-a_3).
	\end{align*}
	
	By property (CH1), we have
	\begin{equation*}
		T_{a_1} \in \mathcal{L_{DP}^{\mathrm{(CH1)}}}(E_2,E_3;F), \quad
		T_{a_2} \in \mathcal{L_{DP}^{\mathrm{(CH1)}}}(E_1,E_3;F), \quad
		T_{a_3} \in \mathcal{L_{DP}^{\mathrm{(CH1)}}}(E_1,E_2;F),
	\end{equation*}
	and hence
	\begin{equation*}
		T_{a_1,a_2} \in \mathcal{DP}(E_3;F), \quad
		T_{a_1,a_3} \in \mathcal{DP}(E_2;F), \quad
		T_{a_2,a_3} \in \mathcal{DP}(E_1;F).
	\end{equation*}
	
	It follows that each term in the above decomposition converges to zero as $j \to \infty$, and therefore
	\begin{equation*}
		\|T(x_j^{(1)},x_j^{(2)},x_j^{(3)}) - T(a_1,a_2,a_3)\| \to 0.
	\end{equation*}
	Thus, $T \in \mathcal{L_{DP}^{\mathrm{ev}}}(E_1,\dots,E_m;F)$.
\end{proof}

The next result provides a condition under which every multilinear operator between Banach spaces is Dunford–Pettis at every point.

\begin{theorem}\label{DPEVS}
	Let $E_1,\dots,E_m$ and $F$ be Banach spaces. If each $E_i$ has the Schur property for $i=1,\dots,m$, then every continuous multilinear operator from $E_1 \times \cdots \times E_m$ into $F$ is Dunford–Pettis at every point. In particular,
	\begin{equation*}
		\mathcal{L_{DP}^{\mathrm{ev}}}(E_1,\dots,E_m;F) = \mathcal{L}(E_1,\dots,E_m;F).
	\end{equation*}
\end{theorem}

\begin{proof}
	Let $E_1,\dots,E_m$ and $F$ be Banach spaces, and assume that each $E_i$ has the Schur property for $i=1,\dots,m$. Let $T \in \mathcal{L}(E_1,\dots,E_m;F)$. We show that $T \in \mathcal{L_{DP}^{\mathrm{ev}}}(E_1,\dots,E_m;F)$.
	
	Let $(x_j^{(i)})_{j=1}^{\infty} \subset E_i$ be such that
	\begin{equation*}
		x_j^{(i)} \overset{w}{\to} a_i \in E_i,
	\end{equation*}
	for all $i=1,\dots,m$. Then
	\begin{equation*}
		x_j^{(i)} - a_i \overset{w}{\to} 0,
	\end{equation*}
	and since each $E_i$ has the Schur property, it follows that
	\begin{equation*}
		x_j^{(i)} - a_i \longrightarrow 0
	\end{equation*}
	in norm, for all $i=1,\dots,m$.
	
	For simplicity, we consider the case $m=3$; the general case follows analogously. Then
	\begin{align*}
		&\|T(x_j^{(1)},x_j^{(2)},x_j^{(3)}) - T(a_1,a_2,a_3)\| \\
		&\le \|T(x_j^{(1)}-a_1, x_j^{(2)}-a_2, x_j^{(3)}-a_3)\| \\
		&\quad + \|T_{a_2,a_3}(x_j^{(1)}-a_1)\| + \|T_{a_1,a_3}(x_j^{(2)}-a_2)\| + \|T_{a_1,a_2}(x_j^{(3)}-a_3)\| \\
		&\quad + \|T_{a_3}(x_j^{(1)}-a_1, x_j^{(2)}-a_2)\| + \|T_{a_2}(x_j^{(1)}-a_1, x_j^{(3)}-a_3)\| \\
		&\quad + \|T_{a_1}(x_j^{(2)}-a_2, x_j^{(3)}-a_3)\| \\
		&\le \|T\| \prod_{i=1}^{3} \|x_j^{(i)} - a_i\|
		+ \|T_{a_2,a_3}\| \|x_j^{(1)} - a_1\|
		+ \|T_{a_1,a_3}\| \|x_j^{(2)} - a_2\| \\
		&\quad + \|T_{a_1,a_2}\| \|x_j^{(3)} - a_3\|
		+ \|T_{a_3}\| \|x_j^{(1)} - a_1\| \|x_j^{(2)} - a_2\| \\
		&\quad + \|T_{a_2}\| \|x_j^{(1)} - a_1\| \|x_j^{(3)} - a_3\|
		+ \|T_{a_1}\| \|x_j^{(2)} - a_2\| \|x_j^{(3)} - a_3\|.
	\end{align*}
	
	Since $x_j^{(i)} - a_i \to 0$ in norm for all $i$, each term on the right-hand side converges to zero as $j \to \infty$. Hence,
	\begin{equation*}
		\|T(x_j^{(1)},x_j^{(2)},x_j^{(3)}) - T(a_1,a_2,a_3)\| \to 0,
	\end{equation*}
	and therefore $T \in \mathcal{L_{DP}^{\mathrm{ev}}}(E_1,\dots,E_m;F)$.
\end{proof}

\begin{corollary}
	Let $E_1,\dots,E_m$ and $F$ be Banach spaces. If each $E_i$ has the Schur property for $i=1,\dots,m$, then
	\[
	\mathcal{L_{DP}}(E_1,\dots,E_m;F) = \mathcal{L_{DP}^{\mathrm{ev}}}(E_1,\dots,E_m;F).
	\]
\end{corollary}

Let us recall that, given Banach spaces $E_1,\dots,E_m$, we denote their projective tensor product by $E_1 \hat{\otimes}_{\pi} \cdots \hat{\otimes}_{\pi} E_m$ (see \cite{R02} for details). As noted in \cite[Section 4]{BT15}, for any Banach spaces $E_1,\dots,E_m$ and $F$, if $E_1 \hat{\otimes}_{\pi} \cdots \hat{\otimes}_{\pi} E_m$ has the Schur property, then
\begin{equation*}
	\mathcal{DP} \circ \mathcal{L}(E_1,\dots,E_m;F) = \mathcal{L}(E_1,\dots,E_m;F).
\end{equation*}
Therefore, we readily obtain the following result.

\begin{proposition}
	Let $E_1,\dots,E_m$ and $F$ be Banach spaces. If each $E_i$ has the Schur property for $i=1,\dots,m$ and $E_1 \hat{\otimes}_{\pi} \cdots \hat{\otimes}_{\pi} E_m$ also has the Schur property, then
	\begin{equation*}
		\mathcal{DP} \circ \mathcal{L}(E_1,\dots,E_m;F)
		= \mathcal{L_{DP}^{\mathrm{ev}}}(E_1,\dots,E_m;F).
	\end{equation*}
\end{proposition}

\begin{proof}
	Since each $E_i$ has the Schur property for $i=1,\dots,m$ and $E_1 \hat{\otimes}_{\pi} \cdots \hat{\otimes}_{\pi} E_m$ also has the Schur property, it follows that
	\begin{equation*}
		\mathcal{L_{DP}^{\mathrm{ev}}}(E_1,\dots,E_m;F)
		= \mathcal{L}(E_1,\dots,E_m;F)
		= \mathcal{DP} \circ \mathcal{L}(E_1,\dots,E_m;F).
	\end{equation*}
\end{proof}


We say that an $m$-linear operator $T \in \mathcal{L}(E_1,\dots, E_m; F)$ is compact if 
$T(B_{E_1}\times \cdots \times B_{E_m})$ is a relatively compact subset of $F$. 
We denote by $\mathcal{K}$ the ideal of compact linear operators and by $\mathcal{L}_K$ 
the ideal of compact multilinear operators. By \cite{P63}, we have
\begin{equation*}
	\mathcal{L}_K = \mathcal{K} \circ \mathcal{L}.
\end{equation*}

Let $E_1,\dots, E_m$ and $F$ be Banach spaces. Assume that each $E_i$, $i=1,\dots,m$, 
and the projective tensor product $E_1 \hat{\otimes}_{\pi} \cdots \hat{\otimes}_{\pi} E_m$ 
have the Schur property. Then, for every $T \in \mathcal{L}(E_1,\dots, E_m; F)$,
\begin{align*}
	T \in \mathcal{L}_K(E_1,\dots, E_m; F) 
	&\iff T_L \in \mathcal{K}(E_1 \hat{\otimes}_{\pi} \cdots \hat{\otimes}_{\pi} E_m; F) \\
	&\subset \mathcal{DP}(E_1 \hat{\otimes}_{\pi} \cdots \hat{\otimes}_{\pi} E_m; F) \\
	&\iff T \in \mathcal{DP} \circ \mathcal{L} 
	= \mathcal{L_{DP}^{\text{ev}}}(E_1,\dots, E_m; F),
\end{align*}
where $T_L$ denotes the linearization of $T$ (see \cite{BPR07} for details). Therefore,
\begin{equation*}
	\mathcal{L}_K(E_1,\dots, E_m; F) 
	\subset \mathcal{L_{DP}^{\text{ev}}}(E_1,\dots, E_m; F).
\end{equation*}

In general, this inclusion is strict, as the following example shows.

\begin{example}
	Let $T\colon \ell_1 \times \cdots \times \ell_1 \longrightarrow \ell_1$ be defined by
	\begin{equation*}
		T\big((x_j^{(1)})_{j=1}^{\infty},\dots,(x_j^{(m)})_{j=1}^{\infty}\big)
		= (x_j^{(1)}\cdots x_j^{(m)})_{j=1}^{\infty}.
	\end{equation*}
	
	Note that $(e_j)_{j=1}^{\infty}$ is bounded in $\ell_1$, since $\|e_j\|_1 = 1$ for all $j$. 
	Moreover,
	\[
	T(e_j,\dots,e_j) = e_j,
	\]
	which has no convergent subsequence. Hence, $T \notin \mathcal{L}_K(\ell_1,\dots,\ell_1; \ell_1)$.
	
	On the other hand, let $(x_j^{(i)})_{j=1}^{\infty}$ be weakly convergent to $x_i$ in $\ell_1$, 
	for $i=1,\dots,m$. Since $\ell_1$ has the Schur property, the convergence is actually 
	in norm. For simplicity, consider $m=3$. Then
	\begin{align*}
		&\|T(x_j^{(1)},x_j^{(2)},x_j^{(3)}) - T(x_1,x_2,x_3)\|_1 \\
		&\le \|x_j^{(1)} - x_1\|_1 \|x_j^{(2)}\|_1 \|x_j^{(3)}\|_1 \\
		&\quad + \|x_1\|_1 \|x_j^{(2)} - x_2\|_1 \|x_j^{(3)}\|_1 \\
		&\quad + \|x_1\|_1 \|x_2\|_1 \|x_j^{(3)} - x_3\|_1.
	\end{align*}
	Since all sequences are bounded and converge in norm, it follows that
	\[
	\|T(x_j^{(1)},x_j^{(2)},x_j^{(3)}) - T(x_1,x_2,x_3)\|_1 \to 0.
	\]
	
	Therefore, $T \in \mathcal{L_{DP}^{\text{ev}}}(\ell_1,\dots,\ell_1; \ell_1)$.
\end{example}


\section{Weakly Dunford-Pettis Multilinear Operators}

We begin by defining the linear case, which will be useful for showing that the class of weakly Dunford-Pettis linear operators is not a hyper-ideal (see \cite{BT15}) and that the coherence and compatibility condition (CH1) fails (see \cite{PR14}).

\begin{definition}
	Let $E$ and $F$ be Banach spaces and let $T \in \mathcal{L}(E;F)$. We say that $T$ is weakly Dunford-Pettis if, for every pair of sequences $(x_j)_{j=1}^{\infty} \subset E$ and $(f_j)_{j=1}^{\infty} \subset F'$ such that $x_j \overset{w}{\to} 0$ and $f_j \overset{w}{\to} 0$, we have
	\begin{equation*}
		f_j\big(T(x_j)\big) \longrightarrow 0.
	\end{equation*}
\end{definition}

We denote by $\mathcal{WDP}(E;F)$ the class of all weakly Dunford-Pettis linear operators from $E$ into $F$. It is straightforward to verify that $(\mathcal{WDP},\|\cdot\|)$ is a Banach operator ideal.

%We will say that a Banach space has the Dunford-Pettis property if the identity, $Id_E\colon E \rightarrow E$ given by $Id_X(x) = x$ for all $x \in E$, is a weakly Dunford-Pettis application.

\begin{definition}
	Let $E_1,\dots,E_m$ and $F$ be Banach spaces, and let $T \in \mathcal{L}(E_1,\dots,E_m;F)$. We say that $T$ is a weakly Dunford-Pettis $m$-linear operator if, for every family of sequences $(x_j^{(i)})_{j=1}^{\infty} \subset E_i$, $i=1,\dots,m$, and every sequence $(f_j)_{j=1}^{\infty} \subset F'$ such that $x_j^{(i)} \overset{w}{\to} 0$ for all $i=1,\dots,m$ and $f_j \overset{w}{\to} 0$, we have
	\begin{equation*}
		f_j\big(T(x_j^{(1)},\dots,x_j^{(m)})\big) \longrightarrow 0.
	\end{equation*}
\end{definition}

We denote by $\mathcal{L_{WDP}}$ the class of all weakly Dunford-Pettis multilinear operators.

\begin{proposition}\label{DPWDP}
	Let $E_1,\dots,E_m$ and $F$ be Banach spaces. Every Dunford--Pettis $m$-linear operator is weakly Dunford--Pettis.
\end{proposition}

\begin{proof}
	Let $T \in \mathcal{L_{DP}}(E_1,\dots,E_m;F)$. Let $(x_j^{(i)})_{j=1}^{\infty} \subset E_i$, $i=1,\dots,m$, and $(f_j)_{j=1}^{\infty} \subset F'$ be such that $x_j^{(i)} \overset{w}{\to} 0$ for all $i=1,\dots,m$ and $f_j \overset{w}{\to} 0$.
	
	Since $T$ is Dunford--Pettis, we have
	\begin{equation*}
		\|T(x_j^{(1)},\dots,x_j^{(m)})\| \longrightarrow 0.
	\end{equation*}
	Moreover, since every weakly convergent sequence is bounded, there exists $K>0$ such that $\|f_j\| \le K$ for all $j$. Therefore,
	\begin{equation*}
		|f_j(T(x_j^{(1)},\dots,x_j^{(m)}))|
		\le \|f_j\| \, \|T(x_j^{(1)},\dots,x_j^{(m)})\|
		\le K \, \|T(x_j^{(1)},\dots,x_j^{(m)})\|.
	\end{equation*}
	Hence,
	\begin{equation*}
		f_j(T(x_j^{(1)},\dots,x_j^{(m)})) \longrightarrow 0,
	\end{equation*}
	and thus $T$ is weakly Dunford--Pettis.
\end{proof}


Following the ideas presented by Ribeiro, Santos and Torres in \cite[Proposition 2.8]{RST22}, we can easily see that $(\mathcal{L_{WDP}}, \|\cdot\|)$, endowed with the usual norm, is a Banach multi-ideal that generalizes the Banach operator ideal $\mathcal{WDP}$.

\begin{remark}
	It is important to emphasize that, in general, the inclusion in Proposition \ref{DPWDP} is strict, since there exist weakly Dunford-Pettis multilinear operators that are not Dunford-Pettis, as will be shown below.
\end{remark}
%%%%%%%%%%%%%%% EXEMPLO RESUMIDO E COMENTADO %%%%%%%%%%%%%%%%%%%%
%\begin{example}
%	Consider the mapping $T\colon c_0 \times c_0 \to c_0$ defined by
%	\[
%	T\big((x_j)_{j=1}^{\infty}, (y_j)_{j=1}^{\infty}\big) = (x_j y_j)_{j=1}^{\infty}.
%	\]
%	Since $e_j \overset{w}{\to} 0$ in $c_0$ and
%	\[
%	\|T(e_j,e_j)\|_{\infty} = \|e_j\|_{\infty} = 1,
%	\]
%	it follows that $T$ is not a Dunford-Pettis bilinear mapping. 
%	
%	On the other hand, using the fact that $c_0' = \ell_1$ and that $\ell_1$ has the Schur property, it is straightforward to verify that $T \in \mathcal{L}_{WDP}(c_0,c_0;c_0)$.
%\end{example}

\begin{example}
	Consider the mapping $T\colon c_0 \times c_0 \to c_0$ defined by
	\[
	T\big((x_j)_{j=1}^{\infty}, (y_j)_{j=1}^{\infty}\big) = (x_j y_j)_{j=1}^{\infty}.
	\]
	Since $e_j \overset{w}{\to} 0$ in $c_0$ and
	\[
	\|T(e_j,e_j)\|_{\infty} = \|e_j\|_{\infty} = 1,
	\]
	it follows that $T$ is not a Dunford--Pettis bilinear mapping.
	
	On the other hand, we claim that $T \in \mathcal{L_{WDP}}(c_0,c_0;c_0)$. Indeed, let $(x_j)_{j=1}^{\infty}$ and $(y_j)_{j=1}^{\infty}$ be sequences in $c_0$ such that $x_j \overset{w}{\to} 0$ and $y_j \overset{w}{\to} 0$, and let $(f_j)_{j=1}^{\infty} \subset \ell_1 = c_0'$ be such that $f_j \overset{w}{\to} 0$. Since $\ell_1$ has the Schur property, it follows that $\|f_j\|_1 \to 0$. Moreover, the sequences $(x_j)$ and $(y_j)$ are bounded, hence there exists $C>0$ such that $\|x_j\|_{\infty}, \|y_j\|_{\infty} \le C$ for all $j$. Therefore,
	\[
	|f_j(T(x_j,y_j))|
	\le \|f_j\|_1 \, \|T(x_j,y_j)\|_{\infty}
	\le \|f_j\|_1 \, \|x_j\|_{\infty} \|y_j\|_{\infty}
	\le C^2 \|f_j\|_1 \longrightarrow 0.
	\]
	Thus, $T$ is weakly Dunford--Pettis.
\end{example}

%Uma vez que todos operadores multilineares Dunford-Pettis é fracamente Dunfurd-Pettis e por Ribeiro, Santos e Torres em \cite[Proposition $2.8$]{RST22} a classe dos operadores  multilineares Dunford-Pettis é um multi-ideal Banach, então, classe dos operadores  multilineares fracamente Dunfurd-Pettis tem que conter os operadores de tipo finito. Vamos então verificar se esta classe é fechada para composição, em outras palavras se esta classe forma um Multi-ideal de operadores multilineares.

%\begin{theorem}\label{WDPMI}
%	A classe $\mathcal{L_{WDP}}$ é um Multi-ideal de aplicações multilineares e quando munido com a norma do operador será um Multi-ideal Banach.
%\end{theorem}
%
%\begin{proof}
%	Sejam $m \in \mathbb{N}$ e $E_1,\dots, E_m, F$ espaços de Banach. Já sabemos que a classe $\mathcal{L_{WDP}}$ contém os operadores de tipo finito, mostraremos então a propriedade de Multi-ideal. Sejam $t \in \mathcal{L}(F; H)$, $T \in \mathcal{L_{WDP}}(E_1,\dots, E_m; F)$ e $\varphi_i \in \mathcal{L}(G_i; E_i)$, $i=1,\dots, m$. Sejam $\left(x_j^{(i)}\right)_{j=1}^{\infty} \subset G_i$ e $(f_j)_{j=1}^{\infty} \subset F'$ tais que $x_j^{(i)} \overset{w}{\longrightarrow}0$ e $f_j\overset{w}{\longrightarrow}0$. Vamos mostrar que $\left(\varphi(x_j^{(i)})\right)_{j=1}^{\infty}$ converge fracamente para zero. Tome $\psi_i \in E_i'$ então
%	\begin{equation*}
%		|\psi_i(\varphi_i(x_j^{(i)}))| = |\psi_i\circ \varphi_i (x_j^{(i)})| \rightarrow 0
%	\end{equation*}
%uma vez que $\psi_i\circ \varphi_i \in G_i'$ e $x_j^{(i)} \overset{w}{\longrightarrow}0$. Dessa forma, $\varphi_i(x_j^{(i)}) \overset{w}{\longrightarrow}0$. Seguindo os mesmos argumentos podemos ver que $t\circ f_j \overset{w}{\longrightarrow}0$. Logo,
%\begin{equation*}
%	f_j\left(t\circ T \circ (\varphi_1,\dots, \varphi_m)(x_j^{(1)},\dots, x_j^{(m)})\right) = (f_j\circ t)\left(T(\varphi_1(x_j^{(1)},\dots, \varphi_m(x_j^{(m)})))\right) \rightarrow 0.
%\end{equation*}
%Dessa forma, $t\circ T \circ (\varphi_1,\dots, \varphi_m) \in \mathcal{L_{WDP}}(G_1,\dots, G_m; H)$.
%
%É imediato ver que, sendo $Id_{\mathbb{K}} \colon \mathbb{K}^m \rightarrow \mathbb{K}$ dada por $Id_{\mathbb{K}}(\lambda_1,\dots, \lambda_m)$ = $\lambda_1\cdots \lambda_m$
%\begin{equation*}
%	\|Id_{\mathbb{K}}\| = 1
%\end{equation*}
%e ainda que, se $t \in \mathcal{L}(F; H)$, $T \in \mathcal{L_{WDP}}(E_1,\dots, E_m; F)$ e $\varphi_i \in \mathcal{L}(G_i; E_i)$, $i=1,\dots, m$, então
%\begin{equation*}
%	\|t\circ T \circ (\varphi_1,\dots, \varphi_m)\| \le \|t\| \|T\| \|\varphi_1\| \cdots \|\varphi_m\|,
%\end{equation*}
%nos mostrando assim que $(\mathcal{L_{WDP}}, \|\cdot\|)$ é um Multi-ideal normado de Aplicações Multilineares. 
%
%Vamos neste momento investigar sua completude em relação à norma do operador. Seja $(T_n)_{n=1}^{\infty}$ uma sequência de Cauchy em $\mathcal{L_{WDP}}(E_1,\dots, E_m; F)$, isto é, para todo $\epsilon > 0$ existe $n_0\in \mathbb{N}$ tal que
%\begin{equation*}
%	\forall m, n \ge n_0 \Rightarrow \|T_n-T_m\| <  \epsilon.
%\end{equation*}
%Uma vez que $\mathcal{L_{WDP}}(E_1,\dots, E_m; F) \subset \mathcal{L}(E_1,\dots, E_m; F)$ então $(T_n)_{n=1}^{\infty}$ também é uma sequência de Cauchy em $\mathcal{L}(E_1,\dots, E_m; F)$ e assim existe $T \in \mathcal{L}(E_1,\dots, E_m; F)$ tal que
%\begin{equation*}
%	T_n \longrightarrow T.
%\end{equation*}
%Vamos mostrar que $T \in \mathcal{L_{WDP}}(E_1,\dots, E_m; F)$. Sejam $\left(x_j^{(i)}\right)_{j=1}^{\infty} \subset E_i$ e $(f_j)_{j=1}^{\infty} \subset F'$ tais que $x_j^{(i)} \overset{w}{\longrightarrow}0$ e $f_j\overset{w}{\longrightarrow}0$
%\begin{equation*}
%	f_j \left(T(x_j^{(1)},\dots, x_j^{(m)})\right) = f_j\left(\lim_{n\rightarrow \infty}T_n (x_j^{(1)},\dots, x_j^{(m)})\right)	= \lim_{n\rightarrow \infty} f_j\left(T_n(x_j^{(1)},\dots, x_j^{(m)})\right) \rightarrow 0.
%\end{equation*}
%Daí, $T \in \mathcal{L_{WDP}}(E_1,\dots, E_m; F)$.
%
%
%\end{proof}

As observed above, the class $(\mathcal{L_{WDP}}, \|\cdot\|)$ is a Banach multi-ideal and therefore enjoys the usual properties of such ideals. It is natural to ask whether this class also satisfies the hyper-ideal property introduced by Botelho and Torres in \cite{BT15}, that is, whether $(\mathcal{L_{WDP}}, \|\cdot\|)$ forms a hyper-ideal of multilinear mappings.

\begin{remark}
	Although $(\mathcal{L_{WDP}}, \|\cdot\|)$ is a Banach multi-ideal, it is not a hyper-ideal of multilinear mappings. Indeed, consider the mapping
	\[
	T\colon \ell_2 \times \ell_2 \to \ell_1, \quad
	T\big((x_j)_{j=1}^{\infty}, (y_j)_{j=1}^{\infty}\big) = (x_j y_j)_{j=1}^{\infty},
	\]
	as in Example 2.9 of \cite{RST22}. It is straightforward to verify that $T \notin \mathcal{L_{WDP}}(\ell_2,\ell_2;\ell_1)$. On the other hand, by Schur's theorem (see \cite[Theorem 1.7]{DJT95}) and the fact that
	\[
	\mathcal{CC}(\ell_1) = \mathcal{DP}(\ell_1) \subset \mathcal{WDP}(\ell_1),
	\]
	we have $Id_{\ell_1} \in \mathcal{WDP}(\ell_1)$. If $(\mathcal{L_{WDP}}, \|\cdot\|)$ were a hyper-ideal, then
	\[
	T = Id_{\ell_1} \circ T \in \mathcal{L_{WDP}}(\ell_2,\ell_2;\ell_1),
	\]
	which is a contradiction.
\end{remark}

\begin{definition}
	Let $m \in \mathbb{N}$, and let $E$ and $F$ be Banach spaces. A mapping $P \in \mathcal{P}(^mE;F)$ is said to be a weakly Dunford-Pettis $m$-homogeneous polynomial if, for every sequence $(x_j)_{j=1}^{\infty} \subset E$ and every sequence $(f_j)_{j=1}^{\infty} \subset F'$ such that $x_j \overset{w}{\to} 0$ and $f_j \overset{w}{\to} 0$, we have
	\begin{equation*}
		f_j\big(P(x_j)\big) \longrightarrow 0.
	\end{equation*}
\end{definition}

%\color{red}
%\begin{proposition}
%	Let $m \in \mathbb{N}$, and let $E$ and $F$ be Banach spaces. For $P \in \mathcal{P}(^mE;F)$, the following are equivalent:
%	\begin{itemize}
%		\item[(i)] $P$ is a weakly Dunford-Pettis $m$-homogeneous polynomial;
%		\item[(ii)] $\check{P} \in \mathcal{L}_{WDP}(^mE;F)$.
%	\end{itemize}
%\end{proposition}
%
%\begin{proof}
%	(i) $\Rightarrow$ (ii) Let $(x_j^{(i)})_{j=1}^{\infty} \subset E$, $i=1,\dots,m$, and $(f_j)_{j=1}^{\infty} \subset F'$ be such that $x_j^{(i)} \overset{w}{\to} 0$ for all $i=1,\dots,m$ and $f_j \overset{w}{\to} 0$. 
%	
%	Using the polarization formula, we can write $\check{P}(x_j^{(1)},\dots,x_j^{(m)})$ as a finite linear combination of terms of the form $P(\varepsilon_1 x_j^{(1)} + \cdots + \varepsilon_m x_j^{(m)})$, where $\varepsilon_k = \pm 1$. Since each sequence
%	\[
%	\varepsilon_1 x_j^{(1)} + \cdots + \varepsilon_m x_j^{(m)} \overset{w}{\to} 0,
%	\]
%	and $P$ is weakly Dunford--Pettis, it follows that
%	\[
%	f_j\big(P(\varepsilon_1 x_j^{(1)} + \cdots + \varepsilon_m x_j^{(m)})\big) \longrightarrow 0.
%	\]
%	Hence, by linearity, we obtain
%	\[
%	f_j\big(\check{P}(x_j^{(1)},\dots,x_j^{(m)})\big) \longrightarrow 0,
%	\]
%	and thus $\check{P} \in \mathcal{L}_{WDP}(^mE;F)$.
%	
%	(ii) $\Rightarrow$ (i) Let $(x_j)_{j=1}^{\infty} \subset E$ and $(f_j)_{j=1}^{\infty} \subset F'$ be such that $x_j \overset{w}{\to} 0$ and $f_j \overset{w}{\to} 0$. Then
%	\[
%	P(x_j) = \check{P}(x_j,\dots,x_j),
%	\]
%	and since $\check{P}$ is weakly Dunford--Pettis, we conclude that
%	\[
%	f_j(P(x_j)) = f_j\big(\check{P}(x_j,\dots,x_j)\big) \longrightarrow 0.
%	\]
%	Therefore, $P$ is weakly Dunford--Pettis.
%\end{proof}
%
%\color{black}

We denote by $\mathcal{P_{WDP}}$ the class of all weakly Dunford-Pettis homogeneous polynomials.

It is possible to relate the class of weakly Dunford--Pettis multilinear mappings to that of weakly Dunford-Pettis homogeneous polynomials through the following result, whose proof follows the same ideas as the proof of condition $(CH5)$ in Theorem \ref{DPcoercomp}.

Before proving this result, we recall an important notion in multilinear theory. Given a Banach multi-ideal $(\mathcal{M}, \|\cdot\|_{\mathcal{M}})$, the associated Banach ideal of homogeneous polynomials is defined by
\begin{equation*}
	\mathcal{P}_{\mathcal{M}} = \{P \in \mathcal{P} \; ; \; \check{P} \in \mathcal{M}\},
\end{equation*}
endowed with the norm
\begin{equation*}
	\|P\|_{\mathcal{P}_{\mathcal{M}}} := \|\check{P}\|_{\mathcal{M}}.
\end{equation*}

\begin{proposition}\label{CHSCHCO}
	Let $m \in \mathbb{N}$, and let $E$ and $F$ be Banach spaces. Then $P \in \mathcal{P_{WDP}}(^mE;F)$ if and only if $\check{P} \in \mathcal{L_{WDP}}(^mE;F)$.
\end{proposition}


%\begin{proof}
%	Sejam $m \in \mathbb{N}$, $E$ e $F$ espaços de Banach. Suponha que $\check{P} \in \mathcal{L_{WDP}}(^mE; F)$, vamo mostrar que $P \in \mathcal{P_{WDP}}(^mE; F)$. Sejam $(x_j)_{j=1}^{\infty} \subset E$ e $(f_j)_{j=1}^{\infty} \subset F'$ tais que $	x_j^{(i)} \overset{w}{\longrightarrow} 0$ e $f_j \overset{w}{\longrightarrow} 0$, asim
%	\begin{equation*}
%		|f_j\left(P(x_j)\right)| = |f_j\left(\check{P}(x_j,\dots, x_j)\right)| \longrightarrow 0.
%	\end{equation*} 
%Suponde agora que $P \in \mathcal{P_{WDP}}(^mE; F)$, considere $(x_j^{(i)})_{j=1}^{\infty} \subset E$ e $(f_j)_{j=1}^{\infty} \subset F'$ tais que $	x_j^{(i)} \overset{w}{\longrightarrow} 0$ e $f_j \overset{w}{\longrightarrow} 0$, para todo $i=1,\dots, m$. Note que, sendo $\epsilon_i = \pm 1$ então
%\begin{equation*}
%	\epsilon_1 x_j^{(1)} + \cdots + \epsilon_m x_j^{(m)} \overset{w}{\longrightarrow} 0
%\end{equation*}
%e assim
%\begin{align*}
%	|f_j\left(\check{P}\left(x_j^{(1)},\dots, x_j^{(m)}\right)\right)| &= \left|f_j\left(\frac{1}{2^mm!}\check{P}(\epsilon_1x_j^{(1)} + \cdots + \epsilon_mx_j^{(m)})^m\right)\right|\\
%	&= \frac{1}{2^mm!}\left|f_j\left(\check{P}(\epsilon_1x_j^{(1)} + \cdots + \epsilon_mx_j^{(m)})^m\right)\right|\\
%	&= \frac{1}{2^mm!}\left|f_j\left(P(\epsilon_1x_j^{(1)} + \cdots + \epsilon_mx_j^{(m)})\right)\right| \longrightarrow 0.
%\end{align*}
%quando $j \rightarrow \infty$. Portanto, $\check{P} \in \mathcal{L_{WDP}}(^mE; F)$.
%\end{proof}

As an immediate consequence of Proposition \ref{CHSCHCO}, we have
\begin{equation*}
	\mathcal{P_{WDP}}(^mE;F) = \mathcal{P}_{\mathcal{L_{WDP}}}(^mE;F)
\end{equation*}
for every $m \in \mathbb{N}$ and all Banach spaces $E$ and $F$.

Another immediate consequence of Proposition \ref{CHSCHCO} is that the pair $(\mathcal{P_{WDP}}(^mE;F), \|\cdot\|)$ is a Banach ideal of homogeneous polynomials.

The result in Proposition \ref{CHSCHCO} corresponds to condition $(CH5)$ in the definitions of coherent and compatible, and strongly coherent and compatible pairs (see \cite{PR14, RST22}). This naturally leads to the question of whether the sequence of pairs $\left(\mathcal{L_{WDP}}, \mathcal{P_{WDP}}\right)$ is strongly coherent and strongly compatible with $\mathcal{WDP}$.

\begin{proposition}
	Let $m,n \in \mathbb{N}$, let $E_1,\dots,E_{m+n}$ and $F$ be Banach spaces. If $T \in \mathcal{L_{WDP}}(E_1,\dots,E_m;F)$ and $Q \in \mathcal{L}(E_{m+1},\dots,E_{m+n})$, then
	\[
	QT \in \mathcal{L_{WDP}}(E_1,\dots,E_{m+n};F).
	\]
\end{proposition}

\begin{proof}
	Let $(x_j^{(i)})_{j=1}^{\infty} \subset E_i$, $i=1,\dots,m+n$, and $(f_j)_{j=1}^{\infty} \subset F'$ be such that $x_j^{(i)} \overset{w}{\to} 0$ for all $i=1,\dots,m+n$ and $f_j \overset{w}{\to} 0$.
	
	Since $Q$ is continuous, the sequence $\big(Q(x_j^{(m+1)},\dots,x_j^{(m+n)})\big)_{j=1}^{\infty}$ is bounded. Moreover, since $T \in \mathcal{L_{WDP}}(E_1,\dots,E_m;F)$, we have
	\[
	f_j\big(T(x_j^{(1)},\dots,x_j^{(m)})\big) \longrightarrow 0.
	\]
	Therefore,
	\[
	\big|f_j\big(QT(x_j^{(1)},\dots,x_j^{(m+n)})\big)\big|
	= \big|Q(x_j^{(m+1)},\dots,x_j^{(m+n)})\big|
	\, \big|f_j\big(T(x_j^{(1)},\dots,x_j^{(m)})\big)\big|
	\longrightarrow 0,
	\]
	which shows that $QT \in \mathcal{L_{WDP}}(E_1,\dots,E_{m+n};F)$.
	
	The norm estimate
	\[
	\|QT\| \le \|Q\| \, \|T\|
	\]
	is immediate.
\end{proof}


Using similar arguments, we obtain the following result.

\begin{proposition}\label{CH3CHCO}
	Let $m,n \in \mathbb{N}$ and let $E$ and $F$ be Banach spaces. If $P \in \mathcal{P_{WDP}}(^mE;F)$ and $Q \in \mathcal{P}(^nE)$, then
	\[
	QP \in \mathcal{P_{WDP}}(^{m+n}E;F).
	\]
\end{proposition}

\begin{remark}\label{WDP(CH1)}
	Although the class $\mathcal{L_{WDP}}$ satisfies conditions $(CH3)$ and $(CH5)$, as shown in Propositions \ref{CHSCHCO} and \ref{CH3CHCO}, it does not satisfy condition $(CH1)$.
	
	Indeed, let $E$ be a Banach space and let $a \in E \setminus \{0\}$. By the Hahn--Banach Theorem, there exists $\varphi \in E'$ such that $|\varphi(a)| = 1$. Define $T\colon E \times \ell_2 \to \ell_2$ by
	\[
	T(x,(x_j)_{j=1}^{\infty}) = (\varphi(x)x_j)_{j=1}^{\infty}.
	\]
	It is immediate that $T$ is a continuous bilinear mapping.
	
	Let $x_j \overset{w}{\to} 0$ in $E$, $y_j \overset{w}{\to} 0$ in $\ell_2$, and $f_j \overset{w}{\to} 0$ in $(\ell_2)'$. Then $\varphi(x_j) \to 0$, while the sequences $(y_j)_{j=1}^{\infty}$ and $(f_j)_{j=1}^{\infty}$ are bounded in $\ell_2$ and $(\ell_2)'$, respectively. Therefore,
	\[
	|f_j(T(x_j,y_j))| \longrightarrow 0,
	\]
	which shows that $T \in \mathcal{L_{WDP}}(E,\ell_2;\ell_2)$.
	
	However, $T_a \notin \mathcal{WDP}(\ell_2;\ell_2)$. Indeed, let $(e_j)_{j=1}^{\infty}$ be the canonical basis of $\ell_2$. Since $(\ell_2)'$ is isometrically isomorphic to $\ell_2$, we may identify each $f_j \in (\ell_2)'$ with $e_j$. Writing $e_j = (\delta_{n,j})_{n=1}^{\infty}$, we obtain
	\[
	|f_j(T_a(e_j))| = |f_j(T(a,e_j))|
	= \left|f_j\big((\varphi(a)\delta_{n,j})_{n=1}^{\infty}\big)\right|
	= |\varphi(a)| \|e_j\|_2
	= |\varphi(a)| = 1,
	\]
	for all $j \in \mathbb{N}$.
	
	Thus, $(T_a(e_j))_{j=1}^{\infty}$ does not converge to zero in $\ell_2$, and hence $T_a$ does not satisfy condition $(CH1)$.
\end{remark}


%Diremos que uma aplicação $m$-linear $T \in \mathcal{L}(E_1,\dots, E_m; F)$ é compacta (fracamente compacta, respectivamente) se $T(B_{E_1}\times \cdots \times B_{E_m})$ é um subconjunto de $F$ relativamente compacto (relativamente fracamente compacto). Chamamos $\mathcal{K}$ e $\mathcal{W}$ os ideais das aplicações lineares compactos e fracamente compactos, respectivamente, e de $\mathcal{L_{K}}$ e $\mathcal{L_{W}}$ os ideais de plaicações multilineares compactos e fracamente compactos respectivamente. 

%\begin{theorem}
%	Sejam $E_1,\dots, E_m$ e $F$ espaços de Banach e $T \in \mathcal{L}(E_1,\dots, E_m; F)$. Então, 
%	\begin{description}
%		\item[(a)] $\mathcal{L_{K}} \subset \mathcal{L_{DP}}$.
%		\item[(b)] $T \in \mathcal{L_{DP}}(E_1,\dots, E_m; F)$ se, e somente se, envia conjuntos relativamente fracamente compactos em um conjunto totalmente limitado.
%		\item[(c)] $T \in \mathcal{L_{WDP}}(E_1,\dots, E_m; F)$ se, e somente se, 
%	\end{description}
%\end{theorem}

%\begin{proposition}\label{CDP}
%	Sejam $E_1,\dots, E_m$ e $F$ espaços de Banach e $T \in \mathcal{L}(E_1,\dots, E_m; F)$. Então, $\mathcal{L_{K}}(E_1,\dots, E_m; F) \subset \mathcal{L_{DP}}(E_1,\dots, E_m; F)$.
%\end{proposition}
%
%\begin{proof}
%	Seja $T \in \mathcal{L_{K}}(E_1,\dots, E_m; F)$. Suponhamos que $T \notin \mathcal{L_{DP}}(E_1,\dots, E_m; F)$, assim existem sequências $(x_j^{(i)})_{j=1}^{\infty}$ in $E_i$, $i=1,\dots, m$, tais que $x_j^{(i)} \overset{w}{\rightarrow} 0$ porém 
%	\begin{equation*}
%		\|T(x_j^{(1)},\dots, x_j^{(m)})\| \not\longrightarrow 0,
%	\end{equation*} 
%quando $j \rightarrow \infty$, isto é, existem subsequências $(x_{j_{k}}^{(i)})_{k=1}^{\infty}$ de $(x_j^{(i)})_{j=1}^{\infty}$ tais que existe $\epsilon_0 > 0$ satisfazendo
%\begin{equation*}
%	\|T(x_{j_k}^{(1)},\dots x_{j_k}^{(m)}) - T(x_{j_l}^{(1)},\dots, x_{j_l}^{(m)})\| \ge \epsilon_0,
%\end{equation*}
%para quaisquer naturais $m$ e $n$ distintos. Por simplicidade, chamaremos $y_k^{(i)} = x_{j_k}^{(i)}$, assim
%\begin{equation*}
%	\|T(y_k^{(1)},\dots, y_k^{(m)}) - T(y_l^{(1)},\dots, y_l^{(m)})\| \ge \epsilon_0
%\end{equation*}
%para quaisquer naturais $m$ e $n$ distintos. Como $(y_k^{(i)})_{k=1}^{\infty}$ é fracamente nula em $E_i$ então, em particular, são limitadas e uma vez que $T \in \mathcal{L_{K}}(E_1,\dots, E_m; F)$ então existem subsequências $(y_{k_p}^{(i)})_{p=1}^{\infty}$ de $(y_k^{(i)})_{k=1}^{\infty}$, para todo $i = 1,\dots, m$, tais que $(T(y_{k_p}^{(1)},\dots, y_{k_p}^{(m)}))_{p=1}^{\infty}$ é convergente, em particular, de Cauchy, o que não ocorre.
%\end{proof}
%
%%\textcolor{red}{\begin{proof}
%%	Seja $T \in \mathcal{L_{K}}(E_1,\dots, E_m; F)$, vamos mostrar que $T \in \mathcal{L_{DP}}(E_1,\dots, E_m; F)$. Sejam $x_j^{(i)}\overset{w}{\longrightarrow} 0$ para todo $i=1,\dots, m$, daí $(x_j^{(i)})_{j=1}^{\infty}$ é uma sequência limitada em $E_i$, dessa forma $\left(T(x_j^{(1)},\dots, x_j^{(m)})\right)_{j=1}^{\infty}$ admite subsequência convergente a zero. Se $\left(T(x_j^{(1)},\dots, x_j^{(m)})\right)_{j=1}^{\infty}$ não convergir a zero, então existe uma subsequência
%%	 $\left(T(x_{j_k}^{(1)},\dots, x_{j_k}^{(m)})\right)_{k=1}^{\infty}$ que não admite subsequências convergindo a zero, porém, uma vez que $(x_{j_k}^{(i)})_{k=1}^{\infty}$ com $i=1,\dots, m$ são limitadas, então $\left(T(x_{j_k}^{(1)},\dots, x_{j_k}^{(m)})\right)_{k=1}^{\infty}$ deveria admitir subsequência convergindo a zero, o que não ocorre. Logo, $T \in \mathcal{L_{DP}}(E_1,\dots, E_m; F)$.
%%\end{proof}}
%
%Uma consequência imediata que vem das Proposições \ref{DPWDP} e \ref{CDP} é a seguinte cadeia de inclusão
%\begin{equation*}
%	\mathcal{L_{K}}(E_1,\dots, E_m; F) \subset \mathcal{L_{DP}}(E_1,\dots, E_m; F) \subset \mathcal{L_{WDP}}(E_1,\dots, E_m; F),
%\end{equation*}
%para quaisquer espaços de Banach $E_1,\dots, E_m$ e $F$.
%
%%\begin{corollary}
%%	Sejam $E_1,\dots, E_m$ e $F$ espaços de Banach e $T \in \mathcal{L}(E_1,\dots, E_m; F)$. Então, $\mathcal{L_{K}}(E_1,\dots, E_m; F) \subset \mathcal{L_{DP}}(E_1,\dots, E_m; F) \subset \mathcal{L_{WDP}}(E_1,\dots, E_m; F)$.
%%\end{corollary}
%
%	Não é difícil ver que se $E_1,\dots, E_m$ e $F$ são reflexivos então os conceitos de Aplicações Multilineares Compactos e Fracamente Dunford-Pettis coincidem. Dessa forma, neste caso temos
%	\begin{equation*}
%		\mathcal{L_{K}}(E_1,\dots, E_m; F) = \mathcal{L_{DP}}(E_1,\dots, E_m; F) = \mathcal{L_{WDP}}(E_1,\dots, E_m; F).
%	\end{equation*}
%
%
%	
%%	De fato, considere $(x_j^{(i)})_{j=1}^{\infty}$ sequências limitadas em $E_i$ e suponha que $T \in \mathcal{L_{WDP}}(E_1,\dots, E_m; F)$. Como cada $E_i$ é reflexivo então existem subsequências, que por simplicidade chamaremos também de $(x_j^{(i)})_{j=1}^{\infty}$ convergindo fracamente para zero em $E_i$. Dessa forma, para qualquer $f_j \overset{w}{\longrightarrow} 0$ em $F'$ temos 
%%	\begin{equation*}
%%		f_j\left(T(x_j^{(1)},\dots, x_j^{(m)})\right) \longrightarrow 0,
%%	\end{equation*}
%%e assim, o conjunto $\left\{T(x_j^{(1)},\dots, x_j^{(m)})|\text{ } j \in \mathbb{N}\right\}$ é um conjunto de Dunford-Pettis. Como $F$ é reflexivo então os conceitos de conjunto de Dunford-Pettis e de conjunto Relativamente compacto coincidem, e portanto, $\left\{T(x_j^{(1)},\dots, x_j^{(m)})|\text{ } j \in \mathbb{N}\right\}$ admite subsequência convergente e assim $T$ é uma Aplicação $m$-linear compacta.

%A bounded subset $A$ of a Banach space $E$ is said to be Dunford-Pettis whenever for any sequences $(x_j)_{j=1}^{\infty}$ in $A$ and $(f_j)_{j=1}^{\infty}$ in $E'$, with $f_j \overset{w}{\longrightarrow} 0$, we have $\lim_{j\rightarrow \infty}f_j(x_j) = 0$.

%\begin{theorem}
%	Let $E_1,\dots, E_m$ and $F$ be Banach spaces and $T \in \mathcal{L}(E_1,\dots, E_m; F)$. The following sentences are equivalent:
%	\begin{enumerate}
%		\item[(a)] $T$ is weakly Dunford-Pettis.
%		\item[(b)] $T$ sends weakly compact subsets of $E_i$ $i=1,\dots, m$ to Dunford-Pettis subsets of $F$.
%		\item[(c)] If $S$ is a weakly compact operator from $Y$ to any arbitrary Banach space, then $S\circ T$ is an $m$-linear Dunford-Pettis application.
%	\end{enumerate}
%\end{theorem}
%
%\begin{proof}
%	$"(a) \Rightarrow (b)"$ Suppose that $T \in \mathcal{L_{WDP}}(E_1,\dots, E_m; F)$. We will make this implication by contradiction. Let $W_i$ be weakly compact subsets in $E_i$ and suppose that there exist $f_j \overset{w}{\longrightarrow} 0$ in $F'$, $(x_j^{(i)})_{j=1}^{\infty}$ in $W_i$ and $\epsilon > 0$ such that
%	\begin{equation*}
%		|f_j(T(x_j^{(1)},\dots, x_j^{(m)}))| > \epsilon.
%	\end{equation*}
%	However, since $W_i$ are weakly compact subsets then, up to one subsequence, we can assume that $x_j \overset{w}{\longrightarrow} 0$ and since $T \in \mathcal{L_{WDP}}(E_1,\dots, E_m; F)$, then
%	\begin{equation*}
%		0 < \epsilon < |f_j(T(x_j^{(1)},\dots, x_j^{(m)}))| \longrightarrow 0
%	\end{equation*}
%	which is absurd. In this way, $T$ sends weakly compact subsets of $E_i$ $i=1,\dots, m$ to Dunford-Pettis subsets of $F$.
%	
%	$"(b) \Rightarrow (c) "$ Let $Z$ be any Banach space and $S\colon F \longrightarrow Z$ a weakly compact linear operator. We will show that $S\circ T$ is an $m$-linear Dunford-Pettis application. Let $x_j^{(i)} \overset{w}{\longrightarrow} 0$ in $E_i$ $i=1,\dots, m$. So, the set
%	\begin{equation*}
%		A = \left\{T(x_j^{(1)},\dots, x_j^{(m)})| \text{ } j \in \mathbb{N}\right\}
%	\end{equation*}
%is a Dunford-Pettis set, and since $S$ is weakly compact, then
%\begin{equation*}
%	S(A) = \left\{S\circ T (x_j^{(1)},\dots, x_j^{(m)})|\text{ } j \in \mathbb{N}\right\}
%\end{equation*}
%is a relatively compact subset, and therefore, up to one subsequence, a subsequence we have that
%\begin{equation*}
%	\|S\circ T(x_j^{(1)},\dots, x_j^{(m)})\| \longrightarrow 0,
%\end{equation*}
%showing that $S\circ T \in \mathcal{L_{DP}}(E_1,\dots, E_m; F)$.
%
%
%%Caso contrário, a sequência $(T(x_j^{(1)},\dots, x_j^{(m)}))_{j=1}^{\infty}$ admite subsequência não convergindo a zero, digamos $(T(x_{j_k}^{(1)},\dots, x_{j_k}^{(m)})_{k=1}^{\infty}$, porém, como $x_j^{(i)} \overset{w}{\longrightarrow} 0$ então
%%\begin{equation*}
%%	\left\{T (x_{j_k}^{(1)},\dots, x_{j_k}^{(m)})|\text{ } k \in \mathbb{N}\right\}
%%\end{equation*}
%%é um conjunto de Dunford-Pettis e assim
%%\begin{equation*}
%%	\left\{S\circ T (x_{j_k}^{(1)},\dots, x_{j_k}^{(m)})|\text{ } k \in \mathbb{N}\right\}
%%\end{equation*}
%%é relativamente compacto e assim
%%
%%\begin{equation*}
%%	\|S\circ T(x_{j_k}^{(1)},\dots, x_{j_k}^{(m)})| \longrightarrow 0,
%%\end{equation*}
%%o que é um absurdo.
%
%$"(c) \Rightarrow (a)"$ Let $x_j^{(i)} \overset{w}{\longrightarrow} 0$ in $E_i$ and $f_j \overset{w}{\longrightarrow} 0$ in $F'$. Consider, $S\colon F \longrightarrow c_0$ given by
%\begin{align*}
%	S(y) = (f_1(y), f_2(y),\dots).
%\end{align*}
%Thus, we can see in \cite[Theorem $17.5$]{AB85} $S$ is weakly compact and thus, $S\circ T$ is an $m$-linear Dunford-Pettis application and so
%\begin{equation*}
%	\|S\circ T(x_j^{(1)},\dots, x_j^{(m)})\|_{\infty} \longrightarrow 0
%\end{equation*}
%and as
%\begin{equation*}
%	|f_j(T(x_j^{(1)},\dots, x_j^{(m)}))| \le \|S\circ T(x_j^{(1)},\dots, x_j^{(m)})\|_{\infty}
%\end{equation*}
%then, $T \in \mathcal{L_{WDP}}(E_1,\dots, E_m; F)$.
%\end{proof}

It is worth noting that, by Theorem \ref{DPEVS}, if $E_i$ has the Schur property for all $i=1,\dots,m$, then
\[
\mathcal{L_{DP}^{\mathrm{ev}}}(E_1,\dots,E_m;F)
= \mathcal{L_{DP}}(E_1,\dots,E_m;F)
= \mathcal{L_{WDP}}(E_1,\dots,E_m;F).
\]

%Another condition that guarantees the equality of classes $\mathcal{L_{DP}}$ and $\mathcal{L_{WDP}}$ is presented in the next result.
%\begin{proposition}\label{DP=WDP}
%Let $E_1,\dots, E_m$ and $F$ be Banach spaces and $F$ reflexive, then
%\begin{equation*}
%	\mathcal{L_{DP}}(E_1,\dots, E_m; F) = \mathcal{L_{WDP}}(E_1,\dots, E_m; F).
%\end{equation*}
%\end{proposition}
%
%\begin{proof}
%	We already know that
%	\begin{equation*}
%		\mathcal{L_{DP}}(E_1,\dots, E_m; F) \subset \mathcal{L_{WDP}}(E_1,\dots, E_m; F)
%	\end{equation*}
%We will then show the other inclusion. Let $T \in \mathcal{L_{WDP}}(E_1,\dots, E_m; F)$ and be weakly null sequences $(x_j^{(i)})_{j=1}^{\infty}$ in $E_i$, with $i=1,\dots, m$ such that $\|T(x_j^{(1)},\dots, x_j^{(m)})\| \nrightarrow 0$. Up to one subsequence we can consider that there exists $\epsilon > 0$ such that
%\begin{equation*}
%	\|T(x_j^{(1)},\dots, x_j^{(m)})\| \ge \epsilon.
%\end{equation*}
% By the Hahn-Banach Theorem, there exists a sequence $(f_j)_{j=1}^{\infty}$ in $F'$ with $\|f_j\| = 1$ and $f_j(T(x_j^{(1)},\dots, x_j^{(m)})) = \|T(x_j^{(1)},\dots, x_j^{(m)})\|$. Thus, by the Banach-Alagolu-Bournaki Theorem, there exists a subsequence $(f_{j_k})_{k=1}^{\infty}$ such that
% \begin{equation*}
% 	f_{j_k} \overset{w*}{\rightarrow} f \in F'
% \end{equation*}
%since $F$ is reflexive, then
%\begin{equation*}
%	f_{j_k} \overset{w}{\rightarrow} f \in F'
%\end{equation*}
%So, as $T \in \mathcal{L_{WDP}}(E_1,\dots, E_m; F)$ we have
%\begin{equation*}
%\|T(x_{j_k}^{(1)},\dots, x_{j_k}^{(m)})\| = |f_{j_k}(T(x_{j_k}^{(1)},\dots, x_{j_k}^{(m)}))| \rightarrow 0
%\end{equation*}
%as $k \rightarrow \infty$. What does not occur.
%\end{proof}

Let us now introduce the concept of weakly Dunford–Pettis multilinear operators at every point, following the same ideas as in Section \ref{DPMOHPEv}.

\begin{definition}\label{WDPML}
	Let $m \in \mathbb{N}$, let $E_1,\dots,E_m$ and $F$ be Banach spaces, and let $T \in \mathcal{L}(E_1,\dots,E_m;F)$. We say that $T$ is a weakly Dunford–Pettis $m$-linear operator at $(a_1,\dots,a_m) \in E_1 \times \cdots \times E_m$ if for every sequence $(x_j^{(i)})_{j=1}^{\infty} \subset E_i$, $i=1,\dots,m$, and every sequence $(f_j)_{j=1}^{\infty} \subset F'$ such that
	\[
	x_j^{(i)} \overset{w}{\longrightarrow} a_i \quad \text{for all } i=1,\dots,m,
	\quad \text{and} \quad
	f_j \overset{w}{\longrightarrow} f \in F',
	\]
	we have
	\[
	f_j\big(T(x_j^{(1)},\dots,x_j^{(m)})\big)
	\longrightarrow
	f\big(T(a_1,\dots,a_m)\big)
	\]
	as $j \to \infty$.
\end{definition}

When $T$ is a weakly Dunford–Pettis multilinear operator at every point of $E_1 \times \cdots \times E_m$, we say that $T$ is weakly Dunford–Pettis at every point and denote the corresponding class by $\mathcal{L_{WDP}^{\mathrm{ev}}}$.

\begin{definition}\label{WDPPH}
	Let $m \in \mathbb{N}$, let $E$ and $F$ be Banach spaces, and let $P \in \mathcal{P}(^mE;F)$. We say that $P$ is a weakly Dunford–Pettis $m$-homogeneous polynomial at $a \in E$ if for every sequence $(x_j)_{j=1}^{\infty} \subset E$ and every sequence $(f_j)_{j=1}^{\infty} \subset F'$ such that
	\[
	x_j \overset{w}{\longrightarrow} a
	\quad \text{and} \quad
	f_j \overset{w}{\longrightarrow} f \in F',
	\]
	we have
	\[
	f_j\big(P(x_j)\big) \longrightarrow f\big(P(a)\big)
	\]
	as $j \to \infty$.
\end{definition}


When $P$ is a weakly Dunford–Pettis homogeneous polynomial at every point of $E$, we say that $P$ is weakly Dunford–Pettis at every point, and we denote the class of all such polynomials by $\mathcal{P_{WDP}^{\mathrm{ev}}}$.

We now present a characterization of weakly Dunford-Pettis multilinear operators at every point, analogous to the Dunford-Pettis case.

\begin{proposition}
	Let $E_1,\dots, E_m$ and $F$ be Banach spaces and $T \in \mathcal{L}(E_1,\dots, E_m; F)$. Then, $T \in \mathcal{L_{WDP}^{\text{ev}}}(E_1,\dots, E_m; F)$ if, and only if, for every $(a_1,\dots, a_m) \in E_1\times\cdots\times E_m$, every family of weakly null sequences $(x_j^{(i)})_{j=1}^{\infty} \subset E_i$, $i=1,\dots,m$, and every sequence $(f_j)_{j=1}^{\infty} \subset F'$ with $f_j \overset{w}{\longrightarrow} f$, we have
	\begin{equation*}
		f_j\big(T(x_j^{(1)} + a_1,\dots, x_j^{(m)} + a_m)\big)
		\longrightarrow
		f\big(T(a_1,\dots, a_m)\big).
	\end{equation*}
\end{proposition}

\begin{proof}
	Suppose first that $T \in \mathcal{L_{WDP}^{\text{ev}}}(E_1,\dots, E_m; F)$.
	
	Let $(a_1,\dots,a_m)\in E_1\times\cdots\times E_m$, let $(x_j^{(i)})$ be weakly null sequences in $E_i$, and let $(f_j) \subset F'$ be such that $f_j \overset{w}{\longrightarrow} f$. Then
	\[
	x_j^{(i)} + a_i \overset{w}{\longrightarrow} a_i
	\]
	for all $i=1,\dots,m$. Since $T$ is weakly Dunford-Pettis at every point, it follows that
	\[
	f_j\big(T(x_j^{(1)} + a_1,\dots, x_j^{(m)} + a_m)\big)
	\longrightarrow
	f\big(T(a_1,\dots, a_m)\big).
	\]
	
	Conversely, suppose the condition holds. Let $(x_j^{(i)}) \subset E_i$ and $(f_j) \subset F'$ be such that
	\[
	x_j^{(i)} \overset{w}{\longrightarrow} a_i
	\quad \text{and} \quad
	f_j \overset{w}{\longrightarrow} f.
	\]
	Define $y_j^{(i)} = x_j^{(i)} - a_i$, so that $y_j^{(i)} \overset{w}{\longrightarrow} 0$. Then
	\[
	f_j\big(T(x_j^{(1)},\dots,x_j^{(m)})\big)
	=
	f_j\big(T(y_j^{(1)} + a_1,\dots,y_j^{(m)} + a_m)\big).
	\]
	By hypothesis,
	\[
	f_j\big(T(y_j^{(1)} + a_1,\dots,y_j^{(m)} + a_m)\big)
	\longrightarrow
	f\big(T(a_1,\dots,a_m)\big).
	\]
	Therefore, $T \in \mathcal{L_{WDP}^{\text{ev}}}(E_1,\dots,E_m;F)$.
\end{proof}

%This characterization shows that weakly Dunford-Pettis multilinear operators at every point are completely determined by their behavior on weakly null perturbations.

Following the same line of reasoning as in Proposition \ref{DPCH1}, and denoting by $\mathcal{L_{WDP}^{\mathrm{(CH1)}}}$ the class of all weakly Dunford-Pettis multilinear operators between Banach spaces that satisfy condition $(CH1)$, we obtain the following result.

\begin{proposition}\label{WDPCH1}
	Let $E_1,\dots,E_m$ and $F$ be Banach spaces. Then,
	\[
	\mathcal{L_{WDP}^{\mathrm{(CH1)}}} = \mathcal{L_{WDP}^{\mathrm{ev}}}.
	\]
\end{proposition}

Using arguments analogous to those in Section \ref{DPML}, the following results can be established.

\begin{theorem}
	\begin{itemize}
		\item[(a)] The class $(\mathcal{L_{WDP}^{\mathrm{ev}}}, \|\cdot\|)$ is a symmetric Banach multi-ideal.
		
		\item[(b)] The class $(\mathcal{P_{WDP}^{\mathrm{ev}}}, \|\cdot\|)$ is a Banach ideal of homogeneous polynomials.
		
		\item[(c)] The sequence $(\mathcal{L_{WDP}^{\mathrm{ev,n}}}, \mathcal{P_{WDP}^{\mathrm{ev,n}}})_{n=1}^{\infty}$ is coherent and compatible with $\mathcal{WDP}$.
	\end{itemize}
\end{theorem}

%\begin{remark}
%	\begin{itemize}
%		\item[(a)] A classe $(\mathcal{L_{WDP}^{\text{ev}}}, \|\cdot\|)$ não é um Hiper-ideal, segundo \cite{BT15}, como mostra o seguinte exemplo (\textcolor{red}{apresentar exemplo}).
%		\item[(b)]  A sequência $(\mathcal{L_{WDP}^{\text{ev}}}_n, \mathcal{P_{WDP}^{\text{ev}}}_n)_{n=1}^{\infty}$ não é fortemente coerente, como mostra o seguinte exemplo (\textcolor{red}{apresentar exemplo}).
%	\end{itemize}
%\end{remark}


\begin{proposition}\label{DPEvsubWDPEv}
	Let $m \in \mathbb{N}$, and let $E_1,\dots,E_m$ and $F$ be Banach spaces. Then,
	\[
	\mathcal{L_{DP}^{\mathrm{ev}}}(E_1,\dots,E_m;F)
	\subset
	\mathcal{L_{WDP}^{\mathrm{ev}}}(E_1,\dots,E_m;F).
	\]
\end{proposition}

\begin{proof}
	Let $T \in \mathcal{L_{DP}^{\mathrm{ev}}}(E_1,\dots,E_m;F)$, and let $(x_j^{(i)})_{j=1}^{\infty} \subset E_i$, $i=1,\dots,m$, and $(f_j)_{j=1}^{\infty} \subset F'$ be such that
	\[
	x_j^{(i)} \overset{w}{\longrightarrow} a_i \in E_i
	\quad \text{for all } i=1,\dots,m,
	\quad \text{and} \quad
	f_j \overset{w}{\longrightarrow} f \in F'.
	\]
	Then,
	\begin{align*}
		&\big|f_j\big(T(x_j^{(1)},\dots,x_j^{(m)})\big) - f\big(T(a_1,\dots,a_m)\big)\big| \\
		&\le \big|f_j\big(T(x_j^{(1)},\dots,x_j^{(m)})\big) - f_j\big(T(a_1,\dots,a_m)\big)\big| \\
		&\quad + \big|f_j\big(T(a_1,\dots,a_m)\big) - f\big(T(a_1,\dots,a_m)\big)\big| \\
		&\le \|f_j\| \, \big\|T(x_j^{(1)},\dots,x_j^{(m)}) - T(a_1,\dots,a_m)\big\|
		+ \big|(f_j - f)\big(T(a_1,\dots,a_m)\big)\big|.
	\end{align*}
	
	Since $(f_j)_{j=1}^{\infty}$ is weakly convergent, it is bounded. Moreover, since $T \in \mathcal{L_{DP}^{\mathrm{ev}}}$, we have
	\[
	T(x_j^{(1)},\dots,x_j^{(m)}) \longrightarrow T(a_1,\dots,a_m)
	\]
	in norm. Hence, the first term converges to zero. The second term also converges to zero because $f_j \to f$ in the weak* topology of $F'$.
	
	Therefore,
	\[
	f_j\big(T(x_j^{(1)},\dots,x_j^{(m)})\big)
	\longrightarrow
	f\big(T(a_1,\dots,a_m)\big),
	\]
	which shows that $T \in \mathcal{L_{WDP}^{\mathrm{ev}}}(E_1,\dots,E_m;F)$.
\end{proof}


%It following ideas similar to those in Proposition \ref{DP=WDP}, we can show that.

\begin{theorem}\label{WDP=WDP}
	Let $E_1,\dots,E_m$ and $F$ be Banach spaces. If $F$ is reflexive, then
	\[
	\mathcal{L_{DP}^{\mathrm{ev}}}(E_1,\dots,E_m;F)
	=
	\mathcal{L_{WDP}^{\mathrm{ev}}}(E_1,\dots,E_m;F).
	\]
\end{theorem}

\begin{proof}
	We already know from Proposition \ref{DPEvsubWDPEv} that
	\[
	\mathcal{L_{DP}^{\text{ev}}}(E_1,\dots, E_m; F) \subset \mathcal{L_{WDP}^{\text{ev}}}(E_1,\dots, E_m; F)
	\]
	for any Banach spaces $E_1,\dots, E_m$ and $F$.
	
	To prove the reverse inclusion, let $T \in \mathcal{L_{WDP}^{\text{ev}}}(E_1,\dots, E_m; F)$. Suppose, by contradiction, that there exist $(a_1,\dots, a_m) \in E_1 \times \cdots \times E_m$ and sequences $(x_j^{(i)})_{j=1}^\infty$ in $E_i$ such that
	\[
	x_j^{(i)} \overset{w}{\longrightarrow} a_i \quad \text{for all } i=1,\dots,m,
	\]
	but
	\[
	T(x_j^{(1)},\dots, x_j^{(m)}) \nrightarrow T(a_1,\dots, a_m).
	\]
	Then, passing to a subsequence if necessary, there exists $\varepsilon > 0$ such that
	\[
	\|T(x_j^{(1)},\dots, x_j^{(m)}) - T(a_1,\dots, a_m)\| \ge \varepsilon
	\]
	for all $j \in \mathbb{N}$.
	
	By the Hahn-Banach Theorem, for each $j$ there exists $f_j \in F'$ with $\|f_j\| = 1$ such that
	\[
	f_j\big(T(x_j^{(1)},\dots, x_j^{(m)}) - T(a_1,\dots, a_m)\big)
	= \|T(x_j^{(1)},\dots, x_j^{(m)}) - T(a_1,\dots, a_m)\|.
	\]
	
	Since $(f_j)$ lies in the unit ball of $F'$, by the Banach--Alaoglu Theorem there exist a subsequence $(f_{j_k})$ and $f \in F'$ such that
	\[
	f_{j_k} \overset{w^*}{\longrightarrow} f.
	\]
	As $F$ is reflexive, the weak* and weak topologies on $F'$ coincide, hence
	\[
	f_{j_k} \overset{w}{\longrightarrow} f.
	\]
	
	Therefore,
	\begin{align*}
		\varepsilon 
		&\le \|T(x_{j_k}^{(1)},\dots, x_{j_k}^{(m)}) - T(a_1,\dots, a_m)\| \\
		&= f_{j_k}\big(T(x_{j_k}^{(1)},\dots, x_{j_k}^{(m)}) - T(a_1,\dots, a_m)\big) \\
		&= f_{j_k}\big(T(x_{j_k}^{(1)},\dots, x_{j_k}^{(m)})\big) - f_{j_k}\big(T(a_1,\dots, a_m)\big).
	\end{align*}
	
	Thus,
	\begin{align*}
		\varepsilon 
		&\le \left| f_{j_k}\big(T(x_{j_k}^{(1)},\dots, x_{j_k}^{(m)})\big) - f\big(T(a_1,\dots, a_m)\big) \right| \\
		&\quad + \left| (f - f_{j_k})\big(T(a_1,\dots, a_m)\big) \right|.
	\end{align*}
	
	Since $T \in \mathcal{L_{WDP}^{\text{ev}}}(E_1,\dots, E_m; F)$ and $f_{j_k} \overset{w}{\longrightarrow} f$, the first term converges to $0$, and the second term also converges to $0$. This is a contradiction.
	
	Therefore,
	\[
	\mathcal{L_{WDP}^{\text{ev}}}(E_1,\dots, E_m; F) \subset \mathcal{L_{DP}^{\text{ev}}}(E_1,\dots, E_m; F),
	\]
	and the proof is complete.
\end{proof}

Another immediate consequence of Theorem \ref{DPEVS} is the following result.

\begin{proposition}
	Let $E_1,\dots, E_m$ and $F$ be Banach spaces. If $E_i$ has the Schur property for all $i=1,\dots, m$, then every continuous multilinear mapping from $E_1\times\cdots\times E_m$ into $F$ is weakly Dunford--Pettis at every point. In particular,
	\[
	\mathcal{L_{WDP}^{\text{ev}}}(E_1,\dots, E_m; F) = \mathcal{L}(E_1,\dots, E_m; F).
	\]
\end{proposition}

\begin{proof}
	By Proposition \ref{DPEvsubWDPEv}, we have
	\[
	\mathcal{L_{DP}^{\text{ev}}}(E_1,\dots, E_m; F)
	\subset 
	\mathcal{L_{WDP}^{\text{ev}}}(E_1,\dots, E_m; F)
	\subset 
	\mathcal{L}(E_1,\dots, E_m; F).
	\]
	Since each $E_i$ has the Schur property, Theorem \ref{DPEVS} ensures that
	\[
	\mathcal{L_{DP}^{\text{ev}}}(E_1,\dots, E_m; F)
	=
	\mathcal{L}(E_1,\dots, E_m; F).
	\]
	The conclusion follows immediately.
\end{proof}

\begin{proposition}
	Let $E_1,\dots, E_m$ and $F$ be Banach spaces. If each $E_i$ has the Schur property and the projective tensor product $E_1 \hat{\otimes}_{\pi} \cdots \hat{\otimes}_{\pi} E_m$ also has the Schur property, then
	\[
	\mathcal{DP} \circ \mathcal{L}(E_1,\dots, E_m; F)
	=
	\mathcal{L_{WDP}^{\text{ev}}}(E_1,\dots, E_m; F).
	\]
\end{proposition}
\begin{proof}
	Since $E_i$ has the Schur property for all $i=1,\dots, m$ and $E_1 \hat{\otimes}_{\pi} \cdots \hat{\otimes}_{\pi} E_m$ has also the Schur property, then
	\begin{equation*}
		\mathcal{L_{WDP}^{\text{ev}}}(E_1,\dots, E_m; F) = \mathcal{L}(E_1,\dots, E_m; F) = \mathcal{DP} \circ \mathcal{L}(E_1,\dots, E_m; F).
	\end{equation*}
\end{proof}


%\begin{proof}
%	By the assumptions, it follows from \cite[Section 4]{BT15} that
%	\[
%	\mathcal{DP} \circ \mathcal{L}(E_1,\dots, E_m; F)
%	=
%	\mathcal{L}(E_1,\dots, E_m; F).
%	\]
%	On the other hand, since each $E_i$ has the Schur property, we know that
%	\[
%	\mathcal{L_{WDP}^{\text{ev}}}(E_1,\dots, E_m; F)
%	=
%	\mathcal{L}(E_1,\dots, E_m; F).
%	\]
%	The result follows immediately.
%\end{proof}
%%%%%%%%%%%%%%%%%%%%%%%%%%%%%%%




%%%%%%%%%%%%%%%%%%%%%%%%%%%%%%%


%Let us recall that given Banach spaces $E_1,\dots, E_m$, we denote the projective tensor product of these spaces by $E_1 \hat{\otimes}_{\pi} \cdots \hat{\otimes}_{\pi} E_m$. For more details see \cite{R02}. Another condition that ensure coincidence results will be presented below. By \cite[Proposition $3.5$]{BPR07}, if $E_1,\dots, E_m$ and $F$ are Banach spaces with $Id_F \in \mathcal{WDP}(F; F)$, then
%\begin{equation*}
%	\mathcal{WDP} \circ \mathcal{L}(E_1,\dots, E_m; F) = \mathcal{L}(E_1,\dots, E_m; F).
%\end{equation*}



%A próxima definição, que nos apresentará uma condição para que os Operadores Multilineares Contínuos Dunford-Pettis sejam Fracamente Dunfor-Pettis pode ser consultada em \cite{AB82}.
%
%\begin{definition}
%	Seja $E$ um espaço de Banach. Diremos que $E$ tem a propriedade de Dunford-Pettis sempre que para quaisquer $x_j \overset{w}{\longrightarrow} 0$ em $E$ e $f_j \overset{w}{\longrightarrow} 0$ em $E'$ temos
%	\begin{equation*}
%		f_j(x_j) \overset{|\cdot|}{\longrightarrow} 0.
%	\end{equation*}
%\end{definition}
%
%
%\begin{proposition}
%	Sejam $E_1,\dots, E_m$ e $F$ Espaços de Banach. Se $F$ admitir a propriedade de Dunford-Pettis, então toda Aplicação Multilinear Dunford-Pettis $T \in \mathcal{L_{DP}}(E_1,\dots, E_m; F)$ é Fracamente Dunford-Pettis.
%\end{proposition}
%\begin{proof}
%	Suponhamos que $F$  admita a propriedade de Dunford-Pettis e seja $T \in \mathcal{L_{DP}}(E_1,\dots, E_m; F)$. Considere então $(x_j^{(i)})_{j=1}^{\infty} \subset E_i$ e $(f_j)_{j=1}^{\infty} \subset F'$ tais que $	x_j^{(i)} \overset{w}{\longrightarrow} 0$ e $f_j \overset{w}{\longrightarrow} 0$, para todo $i=1,\dots, m$ e como $T$ é Dunford-Pettis, então
%	\begin{equation*}
%		\|T(x_j^{(1)},\dots, x_j^{(m)})\| \longrightarrow 0,
%	\end{equation*}
%isto é, $\left(T(x_j^{(1)},\dots, x_j^{(m)})\right)_{j=1}^{\infty}$ é uma sequência convergente a zero na norma e como consequência também será fracamente convergente a zero em $F$. Como $F$ tem a propriedade de Dunford-Pettis então
%\begin{equation*}
%	|f_j(T(x_j^{(1)},\dots, x_j^{(m)}))| \longrightarrow 0.
%\end{equation*}
%Dessa forma, $T$ é fracamente Dunford-Pettis.
%\end{proof}
%
%\textcolor{red}{\begin{example}
%	Denotando por $(r_j)_{j=1}^{\infty}$ a sequência das funções Rademacher sobre $[0, 1]$, isto é, $r_j(t) = sign\left(sen(2^j\pi t)\right)$. Considere a aplicação
%	\begin{equation*}
%		S\colon L[0, 1]^m \longrightarrow \ell_{\infty}
%	\end{equation*}
%definida por, $S(f_1,\dots, f_m) = \left(\int_{0}^{1}f_1(t)\cdots f_m(t) r_j^{+}(t)dt\right)_{j=1}^{\infty}$. Note que, $S$ não é uma aplicação $m$-linear Dunford-Pettis, uma vez que, $r_j \overset{w}{\longrightarrow} 0$ sobre $L[0, 1]$ porém
%\begin{equation*}
%	\|S(r_m,\dots, r_m)\|_{\infty} \ge \int_{0}^{1}r_m(t)\cdots r_m(t)r_m^{+}(t)dt = 2^{-m},
%\end{equation*}
%e assim, $S$ não é Dunford-Pettis. Por outro lado, sendo $f_j \overset{w}{\longrightarrow} 0$ sobre $(\ell_{\infty})'$
%\end{example}}



































%Assim como no caso linear e multilinear podemos definir a classe dos multipolinômios fracamente Dunford-Pettis conforme a definição seguinte.
%
%\begin{definition}
%	Sejam $m, n_1,\dots, n_m \in \mathbb{N}$, $E_1,\dots, E_m$ e $F$ espaços de Banach e $P \in \mathcal{MP}(^{n_1}E_1,\dots,^{n_m}E_m; F)$. Diremos que $P$ é um $(n_1,\dots, n_m)$-polinômio homogêneo fracamente Dunford-Pettis sempre que para quaisquer sequências $(x_j^{(i)})_{j=1}^{\infty} \subset E_i$ e $(f_j)_{j=1}^{\infty} \subset F'$ tais que $	x_j^{(i)} \overset{w}{\longrightarrow} 0$ e $f_j \overset{w}{\longrightarrow} 0$, para todo $i=1,\dots, m$,
%	%	\begin{equation*}
%		%		x_j^{(i)} \overset{w}{\rightarrow} 0 \text{ e } f_j \overset{w}{\rightarrow} 0
%		%	\end{equation*}
%	então
%	\begin{equation*}
%		f_j(P(x_j^{(1)},\dots, x_j^{(m)})) \rightarrow 0.
%	\end{equation*}
%\end{definition}
%
%O conjunto de todos os $(n_1,\dots, n_m)$-polinômios homogêneos entre $E_1,\dots, E_m$ e $F$ será denotado por $\mathcal{MP_{WDP}}(^{n_1}E_1,\dots,^{n_m}E_m; F)$.
%
%Seguindo as mesmas ideias do Teorema \eqref{WDPMI} podemos ver que:
%
%\begin{proposition}
%	A classe $\mathcal{MP_{WDP}}$ é um ideal de $(n_1,\dots, n_m)$-polinômios homogêneos Banach.
%\end{proposition}
%
%Um resultado interessante e inesperado é o seguinte.
%
%\textcolor{red}{\begin{proposition}
%	Sejam $n, m \in \mathbb{N}$, $E_1,\dots, E_m, F$ e $G$ Espaços de Banach. Sendo $P \in \mathcal{P}(^nF; G)$ e $T \in \mathcal{WDP}(E_1,\dots, E_m; F)$ temos que $P \circ T \in \mathcal{WDP}(^nE_1,\dots,^nE_m; G)$ e
%	\begin{equation*}
%		\|P\circ T\| \le \|P\| \|T\|^n.
%	\end{equation*}
%\end{proposition}
%\begin{proof}
%	Sejam $n, m \in \mathbb{N}$, $E_1,\dots, E_m, F$ e $G$ Espaços de Banach. Sejam $\left(x_j^{(i)}\right)_{j=1}^{\infty} \subset G_i$ e $(f_j)_{j=1}^{\infty} \subset F'$ tais que $x_j^{(i)} \overset{w}{\longrightarrow}0$ e $f_j\overset{w}{\longrightarrow}0$, então 
%\end{proof}}

\section{Weakly$^{*}$ Dunford-Pettis Multilinear Operators}

We begin by introducing the linear case. As in the previous section, this approach will be useful to show that the class of weakly$^{*}$ Dunford--Pettis multilinear operators does not satisfy the Coherence and Compatibility condition $(CH1)$ (see \cite{PR14}) and also fails to be a hyper-ideal (see \cite{BT15}).

\begin{definition}
	Let $E$ and $F$ be Banach spaces and $T \in \mathcal{L}(E; F)$. We say that $T$ is \emph{weakly$^{*}$ Dunford--Pettis} if for every sequence $(x_j)_{j=1}^\infty \subset E$ and every sequence $(f_j)_{j=1}^\infty \subset F'$ such that
	\[
	x_j \overset{w}{\longrightarrow} 0
	\quad \text{and} \quad
	f_j \overset{w^{*}}{\longrightarrow} 0,
	\]
	we have
	\[
	f_j\big(T(x_j)\big) \longrightarrow 0.
	\]
\end{definition}

We denote by $\mathcal{WDP}^{*}(E; F)$ the class of all weakly$^{*}$ Dunford-Pettis operators from $E$ into $F$. It is straightforward to verify that $(\mathcal{WDP}^{*}, \|\cdot\|)$ is a Banach operator ideal.

\begin{definition}
	Let $E_1,\dots, E_m$ and $F$ be Banach spaces and $T \in \mathcal{L}(E_1,\dots, E_m; F)$. We say that $T$ is a \emph{weakly$^{*}$ Dunford-Pettis $m$-linear operator} if for every sequence $(x_j^{(i)})_{j=1}^{\infty} \subset E_i$, $i=1,\dots, m$, and every sequence $(f_j)_{j=1}^{\infty} \subset F'$ such that
	\[
	x_j^{(i)} \overset{w}{\longrightarrow} 0 \quad \text{for all } i=1,\dots, m,
	\quad \text{and} \quad
	f_j \overset{w^{*}}{\longrightarrow} 0,
	\]
	we have
	\[
	f_j\big(T(x_j^{(1)},\dots, x_j^{(m)})\big) \longrightarrow 0.
	\]
\end{definition}


\begin{definition}
	Let $m \in \mathbb{N}$, and let $E$ and $F$ be Banach spaces. A polynomial $P \in \mathcal{P}(^mE; F)$ is said to be \emph{weakly* Dunford-Pettis} if for every sequences $(x_j)_{j=1}^{\infty} \subset E$ and $(f_j)_{j=1}^{\infty} \subset F'$ such that
	\begin{equation*}
		x_j \overset{w}{\longrightarrow} 0 \quad \text{and} \quad f_j \overset{w^*}{\longrightarrow} 0,
	\end{equation*}
	we have
	\begin{equation*}
		f_j(P(x_j)) \longrightarrow 0.
	\end{equation*}
\end{definition}

We denote by $\mathcal{L_{WDP^{*}}}$ and $\mathcal{P_{WDP^{*}}}$ the classes of weakly* Dunford-Pettis multilinear operators and homogeneous polynomials, respectively.

\begin{proposition}
	Let $E_1,\dots, E_m$ and $F$ be Banach spaces and $T \in \mathcal{L}(E_1,\dots, E_m; F)$. Then, $T \in \mathcal{L_{WDP^*}^{\text{ev}}}(E_1,\dots, E_m; F)$ if, and only if, for every $(a_1,\dots, a_m) \in E_1\times\cdots\times E_m$, every family of weakly null sequences $(x_j^{(i)})_{j=1}^{\infty} \subset E_i$, $i=1,\dots,m$, and every sequence $(f_j)_{j=1}^{\infty} \subset F'$ with $f_j \overset{w^*}{\longrightarrow} f$, we have
	\begin{equation*}
		f_j\big(T(x_j^{(1)} + a_1,\dots, x_j^{(m)} + a_m)\big)
		\longrightarrow
		f\big(T(a_1,\dots, a_m)\big).
	\end{equation*}
\end{proposition}

\begin{proof}
	Suppose first that $T \in \mathcal{L_{WDP^*}^{\text{ev}}}(E_1,\dots, E_m; F)$.
	
	Let $(a_1,\dots,a_m)\in E_1\times\cdots\times E_m$, let $(x_j^{(i)})$ be weakly null sequences in $E_i$, and let $(f_j) \subset F'$ be such that $f_j \overset{w^*}{\longrightarrow} f$. Then
	\[
	x_j^{(i)} + a_i \overset{w}{\longrightarrow} a_i
	\]
	for all $i=1,\dots,m$. Since $T$ is weakly* Dunford-Pettis at every point, it follows that
	\[
	f_j\big(T(x_j^{(1)} + a_1,\dots, x_j^{(m)} + a_m)\big)
	\longrightarrow
	f\big(T(a_1,\dots, a_m)\big).
	\]
	
	Conversely, suppose the condition holds. Let $(x_j^{(i)}) \subset E_i$ and $(f_j) \subset F'$ be such that
	\[
	x_j^{(i)} \overset{w}{\longrightarrow} a_i
	\quad \text{and} \quad
	f_j \overset{w^*}{\longrightarrow} f.
	\]
	Define $y_j^{(i)} = x_j^{(i)} - a_i$, so that $y_j^{(i)} \overset{w}{\longrightarrow} 0$. Then
	\[
	f_j\big(T(x_j^{(1)},\dots,x_j^{(m)})\big)
	=
	f_j\big(T(y_j^{(1)} + a_1,\dots,y_j^{(m)} + a_m)\big).
	\]
	By hypothesis,
	\[
	f_j\big(T(y_j^{(1)} + a_1,\dots,y_j^{(m)} + a_m)\big)
	\longrightarrow
	f\big(T(a_1,\dots,a_m)\big).
	\]
	Therefore, $T \in \mathcal{L_{WDP^*}^{\text{ev}}}(E_1,\dots,E_m;F)$.
\end{proof}

%This characterization shows that, as in the Dunford-Pettis and weakly Dunford-Pettis cases, weakly* Dunford-Pettis multilinear operators at every point are completely determined by their behavior on weakly null perturbations, with the dual convergence replaced by the weak* topology.


It is straightforward to verify that the class $\left(\mathcal{L_{WDP^{*}}}, \|\cdot\|\right)$, endowed with the usual norm, is a Banach multi-ideal. Moreover, following the same ideas as in Proposition \ref{CHSCHCO}, we have
\begin{equation*}
	\mathcal{P_{WDP^{*}}}(^mE; F) = \mathcal{P}_{\mathcal{L_{WDP^{*}}}}(^mE; F),
\end{equation*}
which ensures that $\left(\mathcal{P_{WDP^{*}}}, \|\cdot\|\right)$ is an ideal of homogeneous polynomials.

We now examine the relationship between this class and the ideals of Dunford-Pettis and weakly Dunford-Pettis multilinear operators.

Since every weakly* convergent sequence is bounded, let $T \in \mathcal{L_{DP}}(E_1,\dots, E_m; F)$. For sequences $(x_j^{(i)})_{j=1}^{\infty} \subset E_i$, $i=1,\dots, m$, with $x_j^{(i)} \overset{w}{\longrightarrow} 0$ and $(f_j)_{j=1}^{\infty} \subset F'$ with $f_j \overset{w^{*}}{\longrightarrow} 0$, we have
\begin{equation*}
	|f_j(T(x_j^{(1)},\dots, x_j^{(m)}))| \le \|f_j\| \, \|T(x_j^{(1)},\dots, x_j^{(m)})\|.
\end{equation*}
Since $(f_j)$ is bounded and $T$ is Dunford-Pettis, it follows that
\begin{equation*}
	f_j(T(x_j^{(1)},\dots, x_j^{(m)})) \longrightarrow 0.
\end{equation*}
Thus, $T \in \mathcal{L_{WDP^{*}}}(E_1,\dots, E_m; F)$, and therefore
\begin{equation*}
	\mathcal{L_{DP}} \subset \mathcal{L_{WDP^{*}}}.
\end{equation*}

Now, let $T \in \mathcal{L_{WDP^{*}}}(E_1,\dots, E_m; F)$ and suppose that $x_j^{(i)} \overset{w}{\longrightarrow} 0$ in $E_i$ for all $i=1,\dots, m$, and $f_j \overset{w}{\longrightarrow} 0$ in $F'$. Since weak convergence implies weak* convergence, we have $f_j \overset{w^{*}}{\longrightarrow} 0$. Hence,
\begin{equation*}
	|f_j(T(x_j^{(1)},\dots, x_j^{(m)}))| \longrightarrow 0,
\end{equation*}
which shows that $T \in \mathcal{L_{WDP}}(E_1,\dots, E_m; F)$. Therefore,
\begin{equation*}
	\mathcal{L_{WDP^{*}}} \subset \mathcal{L_{WDP}}.
\end{equation*}


%It follows immediately from Proposition \ref{DP=WDP} that when $E_1,\dots, E_m$ and $F$ are Banach spaces and $F$ is reflexive, we have
%\begin{equation*}
%	\mathcal{L_{DP}}(E_1,\dots, E_m; F) = \mathcal{L_{WDP^{*}}}(E_1,\dots, E_m; F) = \mathcal{L_{WDP}}(E_1,\dots, E_m; F).
%\end{equation*}

%We say that a subset $A$ of a Banach space $E$ is a bounded set if every weak* null sequence $(f_j)_{j=1}^{\infty}$ of $E'$ converges uniformly on $A$, that is, $\lim_{j\rightarrow \infty}\sup_{x \in A} |f_j(x)| = 0$. It is important to note that every relatively compact set is a bounded set, however, in general, the converse is not true, for example the set $\{e_n \in \ell_{\infty}\text{ ; } n \in \mathbb{N}\}$ is a bounded subset of $\ell_{\infty}$ but is not relatively compact.

%\begin{definition}
%	Let $E_1,\dots, E_m$ and $F$ be Banach spaces and $T \in \mathcal{L}(E_1,\dots, E_m; F)$. We say that $T$ is a bounded Multilinear application if $T(B_{E_1}\times \cdots \times B_{E_m})$ is a bounded subset of $F$.
%\end{definition}

%\begin{theorem}
%	Let $E_1,\dots, E_m, F$ be Banach spaces and $T \in \mathcal{L}(E_1,\dots, E_m; F)$. The following statements are equivalent:
%	\begin{enumerate}
%		\item[(a)] $T \in \mathcal{L_{WDP^{*}}}(E_1,\dots, E_m; F)$.
%		\item[(b)] $T$ sends weakly compact subsets of $E_i$, $i=1,\dots, m$, into a bounded subset of $F$.
%		\item[(c)] If $S_i$ is a weakly compact operator on an arbitrary Banach space $Z_i$, $i=1,\dots, m$, then $T\circ (S_1,\dots, S_m)$ is a bounded $m$-linear mapping.
%		\item[(d)] If $S_i$ is a weakly compact operator from $\ell_1$ to $E_i$, $i=1,\dots, m$, then $T\circ (S_1,\dots, S_m)$ is a bounded $m$-linear application.
%	\end{enumerate}
%\end{theorem}
%
%\begin{proof}
%	$"(a) \Rightarrow (b)"$ Suppose that $T \in \mathcal{L_{WDP^{*}}}(E_1,\dots, E_m; F)$. Let $f_j \overset{w^{*}}{\longrightarrow} 0$ in $F'$ and $A_i$ be weakly compact subsets of $E_i$, $i=1,\dots, m$. We will make this implication by contradiction. Suppose that $(f_j)_{j=1}^{\infty}$ does not converge uniformly to zero on $T(A_1\times \cdots \times A_m)$, then there exists a subsequence of $(f_j)_{j=1}^{\infty}$, that we  call similarly, $(x_j^{(i)})_{j=1}^{\infty}$ in $A_i$, $i=1,\dots, m$, and $\epsilon > 0$ such that
%	\begin{equation*}
%		|f_j(T(x_j^{(1)},\dots, x_j^{(m)}))| > \epsilon.
%	\end{equation*}
%for all $n \in \mathbb{N}$. However, since each $A_i$ is a weakly compact subset, then up to the subsequence, we can assume that $x_j \overset{w}{\longrightarrow} 0$ and since $T \in \mathcal{L_{WDP^{*}}}(E_1,\dots, E_m; F)$, then
%	\begin{equation*}
%		0 < \epsilon < |f_j(T(x_j^{(1)},\dots, x_j^{(m)}))| \longrightarrow 0
%	\end{equation*}
%	which is absurd. In this way, $T$ sends weakly compact subsets of $E_i$ $i=1,\dots, m$ to bounded subsets of $F$.
%	
%	$"(b) \Rightarrow (c)"$ Sejam $Z_i$ espaços de Banach e $S_i\colon Z_i \longrightarrow E_i$, $i=1,\dots, m$, operadores fracamente compactos. Então, $S_i(B_{Z_i})$ é um subconjunto de $E_i$ relativamente fracamente compacto, e assim, 
%	\begin{equation*}
%		T\left(S_1(B_{Z_1})\times \cdots \times S_m(B_{Z_m}) \right) = T\circ (S_1,\dots, S_m)(B_{Z_1}\times \cdots \times B_{Z_m})
%	\end{equation*}
%é um subconjunto limitado de $F$. Dessa forma, $T\circ (S_1,\dots, S_m)$ é uma aplicação $m$-linear limitada.
%
%$"(c) \Rightarrow (d)"$ É imediato.
%
%$"(d) \Rightarrow (a)"$ Sejam $(x_j^{(i)})_{j=1}^{\infty}$ sequências fracamente nulas em $E_i$, $i=1,\dots, m$ e $(f_j)_{j=1}^{\infty}$ uma sequência fraco* nula em $F'$. Considere, $S_i\colon \ell_1 \longrightarrow E_i$ dada por
%\begin{equation*}
%	S_i \left((\lambda_j^{(i)})_{j=1}^{\infty}\right) = \sum_{j=1}^{\infty}\lambda_j^{(i)}x_j^{(i)},\text{ } \forall (\lambda_j^{(i)})_{j=1}^{\infty} \in \ell_1.
%\end{equation*}
%Note que cada $S_i$ é uma aplicação linear fracamente compacta e assim $T\circ (S_1,\dots, s_m)$ é uma aplicação $m$-linear limitada. Dessa forma,
%\begin{equation*}
%	T\circ (S_1,\dots, S_m)(B_{\ell_1}\times \cdots \times B_{\ell_1})
%\end{equation*}
%é um conjunto limitado e assim
%\begin{equation*}
%	\left|f_j(T(x_j^{(1)},\dots, x_j^{(m)})) \right| = \left|f_j(T(S_1(e_1),\dots, S_m(e_m))) \right| = \left|f_j(T\circ (S_1,\dots, S_m)(e_1,\dots, e_m) \right| \longrightarrow 0.
%\end{equation*}
%Assim, $T \in \mathcal{L_{WDP^{*}}}(E_1,\dots, E_m; F)$.

%\end{proof}

\begin{definition}\label{WDP*ML}
	Let $m \in \mathbb{N}$, and let $E_1,\dots, E_m$ and $F$ be Banach spaces. A mapping $T \in \mathcal{L}(E_1,\dots, E_m; F)$ is said to be a \emph{weakly* Dunford-Pettis $m$-linear operator at the point} $(a_1,\dots, a_m) \in E_1 \times \cdots \times E_m$ if for every sequences $(x_j^{(i)})_{j=1}^{\infty} \subset E_i$, $i=1,\dots, m$, and $(f_j)_{j=1}^{\infty} \subset F'$ satisfying
	\begin{equation*}
		x_j^{(i)} \overset{w}{\longrightarrow} a_i \quad \text{and} \quad f_j \overset{w^*}{\longrightarrow} f \in F',
	\end{equation*}
	we have
	\begin{equation*}
		f_j\big(T(x_j^{(1)},\dots, x_j^{(m)})\big) \longrightarrow f\big(T(a_1,\dots, a_m)\big)
	\end{equation*}
	as $j \to \infty$.
\end{definition}

When $T$ is a weakly* Dunford-Pettis multilinear operator at every point of $E_1 \times \cdots \times E_m$, we say that $T$ is \emph{weakly* Dunford-Pettis at every point} and we denote the class of all such operators by $\mathcal{L_{WDP^{*}}^{\text{ev}}}$.

Analogously, we define:

\begin{definition}\label{WDP*PH}
	Let $m \in \mathbb{N}$, and let $E$ and $F$ be Banach spaces. A polynomial $P \in \mathcal{P}(^mE; F)$ is said to be a \emph{weakly* Dunford-Pettis $m$-homogeneous polynomial at the point} $a \in E$ if for every sequences $(x_j)_{j=1}^{\infty} \subset E$ and $(f_j)_{j=1}^{\infty} \subset F'$ satisfying
	\begin{equation*}
		x_j \overset{w}{\longrightarrow} a \quad \text{and} \quad f_j \overset{w^*}{\longrightarrow} f \in F',
	\end{equation*}
	we have
	\begin{equation*}
		f_j(P(x_j)) \longrightarrow f(P(a))
	\end{equation*}
	as $j \to \infty$.
\end{definition}

When $P$ is a weakly* Dunford-Pettis homogeneous polynomial at every point of $E$, we say that $P$ is \emph{weakly* Dunford-Pettis at every point}, and we denote the class of all such polynomials by $\mathcal{P_{WDP^{*}}^{\text{ev}}}$.

It is important to note that, since $\mathcal{L_{WDP^{*}}} \subset \mathcal{L_{WDP}}$, the class $\mathcal{L_{WDP^{*}}}$ inherits several properties from $\mathcal{L_{WDP}}$ established in the previous section. However, it is worth emphasizing that, using the same operator as in Example \ref{WDP(CH1)}, one can readily verify that the class $\mathcal{L_{WDP^{*}}}$ does not satisfy condition $(CH1)$ of \cite{PR14}.

Following the same ideas as in the previous sections, we can obtain the following result.


\begin{theorem}
	\begin{itemize}
		\item[(a)] The class $(\mathcal{L_{WDP^{*}}^{\text{ev}}}, \|\cdot\|)$ is a symmetric Banach multi-ideal.
		
		\item[(b)] $\mathcal{L_{WDP^{*}}^{\text{(CH1)}}} = \mathcal{L_{WDP^{*}}^{\text{ev}}}$.
		
		\item[(c)] The class $(\mathcal{P_{WDP^{*}}^{\text{ev}}}, \|\cdot\|)$ is a Banach ideal of homogeneous polynomials.
		
		\item[(d)] The sequence $(\mathcal{L_{WDP^{*}}^{\text{ev,n}}}, \mathcal{P_{WDP^{*}}^{\text{ev,n}}})_{n=1}^{\infty}$ is coherent and compatible with $\mathcal{WDP^{*}}$.
		
		\item[(e)] Let $m,n \in \mathbb{N}$, $E_1,\dots, E_{m+n}$ and $F$ be Banach spaces. If $T \in \mathcal{L_{WDP^{*}}}(E_1,\dots, E_m; F)$ and $Q \in \mathcal{L}(E_{m+1},\dots, E_{m+n})$, then
		\[
		QT \in \mathcal{L_{WDP^{*}}}(E_1,\dots, E_{m+n}; F)
		\]
		and $\|QT\| \le \|Q\|\,\|T\|$.
	\end{itemize}
\end{theorem}

Since weakly and weakly* convergent sequences are, in particular, bounded, and weak convergence implies weak* convergence, we have
\begin{equation*}
	\mathcal{L_{DP}^{\text{ev}}} \subset \mathcal{L_{WDP^{*}}^{\text{ev}}} \subset \mathcal{L_{WDP}^{\text{ev}}}.
\end{equation*}
Moreover, it follows from Theorem \ref{WDP=WDP} that, if $F$ is reflexive, then
\begin{equation*}
	\mathcal{L_{DP}^{\text{ev}}}(E_1,\dots, E_m; F)
	= \mathcal{L_{WDP^{*}}^{\text{ev}}}(E_1,\dots, E_m; F)
	= \mathcal{L_{WDP}^{\text{ev}}}(E_1,\dots, E_m; F).
\end{equation*}

Following the same arguments as in the previous sections, we obtain:

\begin{proposition}
	Let $E_1,\dots, E_m$ and $F$ be Banach spaces. If $E_i$ has the Schur property for all $i=1,\dots, m$, then every continuous multilinear operator from $E_1 \times \cdots \times E_m$ into $F$ is weakly* Dunford-Pettis at every point. In particular,
	\begin{equation*}
		\mathcal{L_{WDP^{*}}^{\text{ev}}}(E_1,\dots, E_m; F) = \mathcal{L}(E_1,\dots, E_m; F).
	\end{equation*}
\end{proposition}

\begin{proof}
	For any Banach spaces $E_1,\dots, E_m$ and $F$, we have
	\begin{equation*}
		\mathcal{L_{DP}^{\text{ev}}} \subset \mathcal{L_{WDP^{*}}^{\text{ev}}} \subset \mathcal{L_{WDP}^{\text{ev}}} \subset \mathcal{L}.
	\end{equation*}
	If each $E_i$ has the Schur property, then by Theorem \ref{DPEVS} we obtain
	\begin{equation*}
		\mathcal{L_{DP}^{\text{ev}}}(E_1,\dots, E_m; F) = \mathcal{L}(E_1,\dots, E_m; F),
	\end{equation*}
	and the conclusion follows.
\end{proof}

As a consequence, we obtain:

\begin{proposition}
	Let $E_1,\dots, E_m$ and $F$ be Banach spaces. If each $E_i$ has the Schur property and $E_1 \hat{\otimes}_{\pi} \cdots \hat{\otimes}_{\pi} E_m$ also has the Schur property, then
	\begin{equation*}
		\mathcal{DP} \circ \mathcal{L}(E_1,\dots, E_m; F)
		= \mathcal{L_{WDP^{*}}^{\text{ev}}}(E_1,\dots, E_m; F).
	\end{equation*}
\end{proposition}

\begin{proof}
	Under the assumptions, we have
	\begin{align*}
		\mathcal{L_{DP}^{\text{ev}}}(E_1,\dots, E_m; F)
		&= \mathcal{L_{WDP^{*}}^{\text{ev}}}(E_1,\dots, E_m; F)\\
		&= \mathcal{L_{WDP}^{\text{ev}}}(E_1,\dots, E_m; F)\\
		&= \mathcal{L}(E_1,\dots, E_m; F)\\
		&= \mathcal{DP} \circ \mathcal{L}(E_1,\dots, E_m; F),
	\end{align*}
	which completes the proof.
\end{proof}


\begin{thebibliography}{C}
	
	\bibitem{AB82} Aliprantis, C. D.; Burkinshaw O.: {\it Dunford-Pettis Operator on Banach Lettices}. Transactions of the american mathematical society, {\bf 274}, (1982).
	
	\bibitem{AB85} Aliprantis, C. D.; Burkinshaw O.: {\it Positive operators}, Reprint of the 1985 original. Springer, Dordrecht, (2006).
	
	%\bibitem{37} Albuquerque N.,  Ara\'ujo G., Pellegrino D., Rueda P., A note multiple summing operators and applications, https://arxiv.org/pdf/1503.07306v3.pdf, (2015).
	
	\bibitem{AD02} Aron R., Dimant V., {\it Sets of weak sequential continuity for polynomials}. Indag. Math., {\bf 13(3)}, 287–299, (2002).
	
	%\bibitem{AR12} Aron, RM.; Rueda, P.: {\it Ideals of homogeneous polynomials}. Publ. Res. Inst. Math. Sci. {\bf 46}, 957–969, (2012).
	
	%\bibitem{AR15} Aron, RM.; Rueda, P.: {\it $\mathcal{I}$-bounded holomorphic functions}. J. Nonlinear Convex Anal. {\bf 16}, 2539–2551, (2015).
	
	%\bibitem {Banach32} Banach, S.: {\it Theorie de operations lineaires}. Chelsea Publishing Company, (1932).
	
	%\bibitem{BPSS15} Bernardino, A.; Pellegrino, D.; Seoane-Sep\'ulveda, J. B.; Souza, M. L. V.: {\it Nonlinear absolutely summing operators revisited}. Sociedade Brasileira de Matematica. Boletim, Nova Serie, {\bf 46}, 205-249, (2015).
	
	%\bibitem{B05} Botelho, G.: {\it Ideals of polynomials generated by weakly compact operators}, Note Mat. {\bf 25}, 69-102, (2006).
	
	%\bibitem{BBDP07} Barbosa, J., Botelho, G., Diniz, D., Pellegrino, D.: {\it Spaces of absolutely summing	polynomials}, Mathematica Scandinavica. {\bf 101}, 219-237, (2007).
	
	%\bibitem{BBJ01} Botelho, G.; Braunss, H. A.; Junek, H.: {\it Almost $p$ - summing polynomials and multilinear mappings}. Archiv der Mathematik. 76, 109-118, (2001).
	
	%\bibitem{BBJP06} Botelho, G.; Braunss, H.-A.; Junek, H.; Pellegrino, D.: {\it Holomorphy types and ideals of multilinear mappings}. Studia Math. {\bf 177}, 43-65, (2006).
	
	%\bibitem{18} Botelho G., H.-A. Braunss, H. Junek and D. Pellegrino, Holomorphy types and ideals of multilinear mappings, Studia Mathematica 177 (2006), 43–65.
	
	%\bibitem{19} Botelho, G.; Braunss H.; Junek H.; Pellegrino, D.: {\it Inclusions and Coincidences for Multiple Summing Multilinear Mappings}. Proceedings of the American Mathematical Society, vol 137 (3), 991-1000, (2009).
	
	\bibitem{BHJ22} Benkhaled, H.; Hajji, M.; Jeribi, A.: {\it On the class of demi Dunford-Pettis operators}. Rendiconti del Circolo Matematico di Palermo. (2022)
	
	\bibitem{BKAB24} Boumnidel S.; Kaddouri, A. El.; Aboutafail, O.; Bouras K.: {\it Some Properties of Weak* Dunford-Pettis Operators on Banach Lattices}. Bol. Soc. Paran. Mat., {\bf 42}, (2024).
	
	
	%\bibitem{BC17} Botelho, G.; Campos, J.: {\it On the transformation of vector-valued sequences by linear and multilinear operators}. Monatsh. Math. {\bf 183}, 415-435, (2017).
	
	%\bibitem{BCS17} Botelho, G.; Campos, J.; Santos J.: {\it Operator ideals related to absolutely summing and Cohen strongly summing operators}. Pac. J. Math. {\bf 287}, 1-17, (2017).
	
	\bibitem{BPR07} Botelho, G.; Pellegrino, D.; Rueda, P.:
	{\it On Composition ideals of multilinear mappings and homogeneous polynomials}. Publications of the Research Institute for Mathematical Sciences.  {\bf 43}, 1139-1155, (2007).
	
	\bibitem{BT15} Botelho, G.;  Torres, E. R.: {\it Hyper-ideals of multilinear operators}, Linear Algebra Appl. {\bf 482}, 1--20, (2015).
	
	%\bibitem{BT17} Botelho, G.; Torres,  E. R.: {\it Techniques to generate hyper-ideals of multilinear operators}. Linear Multilinear Algebra {\bf 65}, 1232--1246, (2017).
	
	\bibitem{BT18} Botelho, G.; Torres, E. R.:{\it Two-sided polynomial ideals on Banach spaces}. Journal of Mathematical Analysis and Applications. {\bf 462}, 900-914, (2018).
	
	
	%\bibitem{BTV18} Botelho, G.; Torres, E.; Velanga, T.: {\it Linearization of multipolynomials and applications}. Arch. Math. {\bf 110}, 605-615, (2018).
	
	%\bibitem{BP05} Botelho, G.; Pellegrino D.: {\it Two new properties of ideals of polynomials and applications}. Indagationes Mathematicae. vol 16 (2), 157-169, (2005).
	
	%\bibitem{20} Botelho, G; Pellegrino, D: {\it When every multilinear mapping is multiple summing}. Mathematische Nachrichten, vol 282, 1414-1422, (2009).
	
	%\bibitem{BW22} Botelho, G.; Wood, R.: {\it Hyper-ideals of Multilinear Operators and Two-Sided Polynomial Ideals Generated by Sequence Classes}. Mediterranean Journal of Mathematics. {\bf 20}, (2022).
	\bibitem{BR98} Boyd C., Ryan R.A.; {\it Bounded weak continuity of homogeneous polynomials at the origin}. Arch. Math. (Basel), {\bf 71}, 211–218, (1998).
	
	%\bibitem{BJ} Braunss, H.-A.; Junek, H.; {\it Ideals of Polnomials and Multilinear Mappings}. Unpublished notes.
	
	%\bibitem{C14} Campos, J.: {\it Cohen and multiple Cohen strongly summing multilinear operators}. Linear Multilinear Algebra. {\bf 62}, 322-346, (2014).
	
	%\bibitem{35} Campos J.:  Contribui\c c\~oes \`a teoria dos operadores Cohen fortemente somantes. Tese de Doutorado em Matemática - Universidade Federal da Paraí­ba, Jo\~ao Pessoa, (2013).
	
	%\bibitem{27} Carando, D.; Dimant, V.; Muro, S.: {\it Coherent sequences of polynomial ideals on Banach spaces}. Mathematische Nachrichten. vol 282, 1111-1133, (2009).
	
	%\bibitem{29} Carando, D.; Dimant, V.; Muro, S.: {\it Every Banach ideal of polynomials is compatible with an operator ideal}. Monatshefte fur Mathematik. vol 165, 1-14, (2012).
	
	%\bibitem{28} Carando, D.; Dimant, V.; Muro, S.: {\it Holomorphic functions and polynomial ideals on Banach spaces}. Collectanea Mathematica. vol 63, 71-91, (2012).
	
	\bibitem{DJT95} Diestel J., Jarshow H., A. Tonge, Absolutely Summing Operators. Cambridge Studies in Advanced Mathematics, Cambridge, (1995).
	
	\bibitem{DU77}  Diestel, J. J.; Uhl, J. J.: {\it Vector measures}. Math. Surveys, No. 15, Amer. Math. Soc, Providence, R. I (1977).
	
	%\bibitem{1} Diestel J., Jarchow H., Pietsch A., Operator Ideals, Handbook of the Geometry of Banach Spaces,
	%Vol 1, Elsevier (2001), 437-496.
	
	%\bibitem{21} Defant, A.; Popa, D.; Schwarting, U.: {\it Coordinatewise multiple summing operators in Banach spaces}. Journal of Functional Analysis. vol 259, 220-242, (2010).
	
	%\bibitem {DR50} Dvoretzky, A. and Rogers, C. A. \begin{itshape} Absolutely and unconditional convergence in normed spaces\end{itshape}. Proc. Nat. Acad. Sci. USA, 36 (1950), 192-197.
	
	\bibitem{FG03} Floret, K.; Garc\'ia, D.: {\it On ideals of polynomials and multilinear mappings between Banach spaces}. Arch. Math. {\bf 81}, 300-308, (2003).
	
	%\bibitem{G04} Garc\'ia, D.: {\it The inclusion theorem for multiple summing operators}. Studia Mathematica. vol 165 (3), 275-290, (2004).
	
	%\bibitem {G55} Grothendieck A., \begin{itshape} Produits tensoriels topologiques et espaces nucl\'eaires\end{itshape}. Memoirs Acad. Math. Soc. 16, (1955).
	
	%\bibitem{2} Grothendieck A., R\`esum\`e de la th\`orie m\`etrique des produits tensoriels topologiques, Boletim da
	%Sociedade Matem\'atica de S\~ao Paulo 8, 1-79, (1956).
	
	%\bibitem{KS14} Karn, A., Sinha, D.: An operator summability of sequences in Banach spaces. Glasgow Mathematical Journal. {\bf 56}, 427-437, (2014).
	
	%\bibitem {LP68} Lindenstrauss J., Pe\l czy\'{n}ski A., \begin{itshape} Absolutely summing operators in $L_{p}$ spaces and their applications\end{itshape}. Studia Math. 29, 276-326, (1968).
	
	%\bibitem{39} Garc\'ia D.; Villanueva, I.: {\it Multiple summing operators on Banach spaces}. Journal of Mathematical Analysis and Applications, vol 285 (1), 86-96, (2003).
	
	%\bibitem{24} Matos, M. C.: {\it Absolutely summing mapping, nuclear mappings and convolution equations}. $http:www.ime.unicamp.br/rel_pesq/2007/pdf/rp03-07.pdf$, (2005).
	
	%\bibitem{26} Matos, M. C.: {\it Fully absolutely summing and Hilbert-Schmidt multilinear mapping}. Collectanea Mathematica. vol 54, 111-136, (2003).
	
	%\bibitem{M03} Matos, M.: {\it Nonlinear absolutely summing mappings}, Math. Nachr. {\bf 258}, 71-89, (2003).
	
	%\bibitem{M04} Matos, M.: {\it Mappings between Banach spaces that send mixed summable sequences into absolutely summable sequences}. Journal of Mathematical Analysis and Applications. {\bf 297}, 833-851, (2004).
	
	
	%\bibitem {SB} Maudlin R. D., \begin{itshape} The Scottish Book, Mathematics from the Scottish Caf\'e\end{itshape}. Birkhauser Verlag, Boston-Basel-Stuttgart (1981).
	
	%\bibitem{M10} Mujica, J.: {\it Complex Analysis in Banach Spaces}. Dover publications, Mineola, New
	%York, (2010).
	
	\bibitem{MN91} Meyer-Nieberg, P.: Banach Lattices. Springer, Berlin (1991)
	
	%\bibitem{PR12} Pellegrino, D.; Ribeiro, J.: {\it On almost summing polynomials and multilinear mappings. Linear and Multilinear Algebra}. Linear Multilinear Algebra. {\bf 60}, 397-413, (2012).
	
	\bibitem{P63} Pelczynski, A.: {\it On weakly compact polynomial operators on B-spaces with Dunford-Pettis Property}. Bulletin L'Académie Polonaise des Science, Série des Sciences Mathé-matiques, Astronomiques et Physiques. {\bf 11}, 371-377, (1963).
	
	\bibitem{PR14} Pellegrino, D.; Ribeiro, J.: {\it On multi-ideals and polynomial ideals of Banach spaces: a new approach to coherence and compatibility}. Monatshefte fur Mathematik. {\bf 173}, 379-415, (2014).
	
	%\bibitem{PR12} Pellegrino, D.; Ribeiro, J.: {\it On almost summing polynomials and multilinear mappings. Linear and Multilinear Algebra}. Linear Multilinear Algebra. {\bf 60}, 397-413, (2012).
	
	%\bibitem{32} Pellegrino, D.; Seoane-Sep\'ulveda J. B.: {\it Multiple $(p, q, r)$-summing polynomial and multilinear operators}. ResearchGate, (2011). PARECE QUE ESSE ARTIGO É O DE NÚMERO 3. PROCUREI E NÃO ENCONTREI ELE NO LATTES DE DANIEL.
	
	%\bibitem{P67} Pietsch, A.: {\it Absolutely $p$-summierende Abbildungen in normieten Raumen}. Studia Math. {\bf 27}, 333-353, (1967).
	
	%\bibitem{P80} Pietsch A.: {\it Operator Ideals}. Deutscher Verlag der Wiss, Berlin (1978) [North Holland, Amsterdam (1980)].
	
	\bibitem{P83} Pietsch A.: {\it Ideals of multilinear functionals}. Proceedings of the second international conference on operator algebras, ideals, and their applications in theoretical physics, 185-199, Teubner, Leipzig (1983).
	
	
	
	%\bibitem {G56} A. Grothendieck, \begin{itshape} Resum\'e de la th\'eorie m\'etrique des produits %tensoriels topologiques\end{itshape}. Bol. Soc. Mat. Sao Paulo. 8 (1956), 1-79.
	
	
	%\bibitem{3} Ramanujan M. S., Absolutely $\lambda$-summing operators,$\lambda$ a symmetric sequence space,Math.Z.114, 187-193, (1970).
	
	%\bibitem{R11} Ribeiro, J., Ideais coerentes e compat\'iveis entre espa\c cos de Banach. Tese de Doutorado em Matem\'atica - Universidade Federal de Pernambuco, (2011).
	
	%\bibitem{RS21} Ribeiro, J., Santos F., {\it Absolutely Summing Polynomials}. Methods of Functional Analysis and Topology, {\bf 27}, 74-87, (2021).
	
	%\bibitem{RS19} Ribeiro, J.; Santos, F.: {\it Generalized Multiple Summing Multilinear Operators on Banach Spaces}, Mediterranean Journal of Mathematics. {\bf 16}, n. 108, (2019).
	
	%\bibitem{RSS21} Ribeiro J., Santos D., Santos F.: {\it Multipolynomials: An Almost Symmetrical Approach}. Results Math, {\bf 76}, 159, (2021).
	
	\bibitem{RST22} Ribeiro J., Santos F., Torres E.R.: {\it Coherence and compatibility: a stronger approach}.
	Linear and multilinear algebra, {\bf 70}, 66-80, (2022).
	
	\bibitem{R02} Ryan R.A.: {\it Introduction to Tensor Products of Banach Spaces}. Springer-Verlag, (2002).
	
	%\bibitem{S20} Santos, F. A. O.: {\it Contribuições à teoria da coerência e
	%	compatibilidade}. Tese de doutorado, (2020).
	
	%\bibitem{S13} Serrano-Rodr\'iguez D. M.: {\it Absolutely $\gamma$-summing multilinear operators}, Linear Algebra Appl. {\bf 439}, 4110-4118, (2013).
	
	%\bibitem{S80} Stephani I.: {\it Generating systems of sets and quotients of surjective operator ideals}. Math.Nachr. {\bf 99}, 13–27, (1980).
	
	%\bibitem{VZ11} Vasylyshyn, T. V.; Zagorodnyuk, A. V.: {\it Polarization formula for $(p, q)$-polynomials on a complex normed space}. Methods Funct. Anal. Topology {\bf 17}, no. 1, 75–83, (2011).
	
	%\bibitem{V18} Velanga, T.: {\it Ideals of polynomials between Banach spaces revisited}, Linear Multilinear Algebra. {\bf 66}, 2328-2348, (2018).
	
	%\bibitem{V19} Velanga, T.; {\it Absolutely summing multipolynomials}, Results Math. {\bf 74} , p. 19, (2019).
	
	
	
	%\bibitem{V17MP} Velanga, T.: {\it Multilinear applications versus homogeneous polynomials}, \textcolor{blue}{arXiv:1706.04703} [math.FA], (2017).
	
	%\bibitem{8} J. Lindenstrauss e A.Pelczy\'nski, Absolutely Summing operators in $L_p$ spaces and their applications, Studia Mathematica 29 (1968), 275-326.
	
	%\bibitem{23} Botelho, G. and Campos, J.R. Monatsh Math (2016). doi:10.1007/s00605-016-0963-4.
	
	%\bibitem \begin{itshape} \end{itshape}
	
	
	
\end{thebibliography}

\end{document}
